\documentclass[11pt,a4paper]{article}

% ---------- Encodage et langue ----------
\usepackage[utf8]{inputenc}
\usepackage[T1]{fontenc}
% \usepackage{lmodern}  % décommenter si disponible (police vectorielle)
\usepackage[french]{babel}
\usepackage[a4paper,margin=2.3cm]{geometry}

% ---------- Mathématiques ----------
\usepackage{amsmath,amssymb,amsthm}
\usepackage{mathtools}
\usepackage{enumitem}
\usepackage[most]{tcolorbox}
\usepackage[colorlinks=true,linkcolor=blue!60!black,urlcolor=blue!60!black]{hyperref}

% ---------- Macros ----------
\newcommand{\R}{\mathbb{R}}
\newcommand{\N}{\mathbb{N}}
\newcommand{\Z}{\mathbb{Z}}
\newcommand{\Q}{\mathbb{Q}}
\newcommand{\C}{\mathbb{C}}
\newcommand{\K}{\mathbb{K}}
\newcommand{\Sun}{\mathbb{S}^1}
\newcommand{\interieur}[1]{\mathring{#1}}
\newcommand{\adh}[1]{\overline{#1}}
\newcommand{\Int}{\operatorname{int}}
\newcommand{\diam}{\operatorname{diam}}
\newcommand{\dist}{\operatorname{dist}}
\newcommand{\Vect}{\operatorname{Vect}}
\newcommand{\Bf}{B_f}
\newcommand{\norm}[1]{\left\|#1\right\|}
\newcommand{\abs}[1]{\left|#1\right|}
\newcommand{\opnorm}[1]{|\!|\!|#1|\!|\!|}
\newcommand{\ps}[2]{\left\langle #1,#2\right\rangle}
\newcommand{\Lc}{\mathcal{L}}
\newcommand{\Cz}{\mathcal{C}^0}
\newcommand{\Cinf}{\mathcal{C}^\infty}
\newcommand{\Cc}{\mathcal{C}^0_c}
\newcommand{\ie}{c'est-à-dire}
\newcommand{\wk}{\rightharpoonup}
\newcommand{\wks}{\rightharpoonup^{\!\star}}
\DeclareMathOperator{\Ker}{Ker}
\DeclareMathOperator{\Ima}{Im}
\DeclareMathOperator{\supp}{supp}
\DeclareMathOperator{\sgn}{sgn}
\renewcommand{\Re}{\operatorname{Re}}
\renewcommand{\Im}{\operatorname{Im}}
\newcommand{\D}{\mathcal{D}}
\newcommand{\Sha}{\mathrm{III}}

% ---------- Environnements ----------
\newtcolorbox{exo}[1]{
  enhanced, breakable,
  colback=black!3, colframe=black!70,
  fonttitle=\bfseries, title={#1},
  boxrule=0.6pt, arc=1pt, left=6pt, right=6pt
}
\newtcolorbox{pointcle}{
  enhanced, breakable,
  colback=blue!4, colframe=blue!55!black,
  fonttitle=\bfseries, title={Point clé},
  boxrule=0.8pt, arc=2pt, left=6pt, right=6pt
}
\newtcolorbox{piege}{
  enhanced, breakable,
  colback=red!4, colframe=red!60!black,
  fonttitle=\bfseries, title={Piège / difficulté},
  boxrule=0.8pt, arc=2pt, left=6pt, right=6pt
}
\newtcolorbox{commentaire}{
  enhanced, breakable,
  colback=gray!8, colframe=gray!55!black,
  fonttitle=\bfseries, title={Commentaire},
  boxrule=0.6pt, arc=2pt, left=6pt, right=6pt
}
\newtcolorbox{methode}{
  enhanced, breakable,
  colback=green!4, colframe=green!40!black,
  fonttitle=\bfseries, title={Méthode},
  boxrule=0.8pt, arc=2pt, left=6pt, right=6pt
}

\theoremstyle{definition}
\newtheorem*{solution}{Solution}

\title{\textbf{4MA305 -- Bases d'analyse fonctionnelle}\\[2mm]
\Large Corrigé détaillé du TD n\textsuperscript{o}6 -- Distributions}
\author{}
\date{Année universitaire 2026--2027}

\begin{document}
\maketitle

\begin{commentaire}
\textbf{Conventions.} $\Omega$ est un ouvert de $\R^d$ ; $\D(\Omega)$ l'espace des fonctions tests, $\D_K(\Omega)$ celles à support dans le compact $K$, $\norm\varphi_{m,K}=\max_{|\alpha|\le m}\norm{\partial^\alpha\varphi}_K$ ; $\D'(\Omega)$ l'espace des distributions (Définition~6.3), avec l'estimation de continuité (6.1), la convergence $T_n\wk T$ de la Définition~6.4, la dérivation $\langle\partial^\alpha T,\varphi\rangle=(-1)^{|\alpha|}\langle T,\partial^\alpha\varphi\rangle$, le produit $\langle fT,\varphi\rangle=\langle T,f\varphi\rangle$ ($f\in\Cinf$), la translation $\langle\tau_aT,\varphi\rangle=\langle T,\tau_{-a}\varphi\rangle$ et la convolution $\langle\psi\star T,\varphi\rangle=\langle T,\varphi\star\check\psi\rangle$, $\check\psi(x)=\psi(-x)$. On identifie $f\in L^1_{\mathrm{loc}}(\Omega)$ à $T_f$ (Lemme~6.2) et on écrit souvent $f$ pour $T_f$, $f'$ pour $T_f'$. $h=\mathbf 1_{\R_+}$ est la fonction de Heaviside, $\delta_a$ la masse de Dirac, $(\rho_n)$ une approximation de l'unité issue d'une fonction test positive paire à support dans la boule unité. Le Théorème~6.4 (distributions à dérivée nulle sur un connexe), la formule des sauts (Théorème~6.3), la Proposition~6.10 (régularité de $\psi\star T$) et le Corollaire~6.2 sont supposés connus.

Cette feuille ne figure pas parmi les feuilles officielles ; elle complète les exercices du polycopié et les développements du TD0 (continuité séquentielle, mesures de Radon, séries trigonométriques, ordre, formule des sauts et peigne de Dirac, valeur principale, $y'=ay$, solution fondamentale du laplacien dans le plan), que l'on cite par « TD0 ch.~6 ». Les encadrés « Point clé » signalent les idées à retenir, « Piège » les erreurs fréquentes, « Méthode » les schémas de preuve réutilisables.
\end{commentaire}

\section*{Fonctions tests, distributions, convergence}

% =====================================================================
\begin{exo}{Exercice 1 -- Fabriquer des fonctions tests}
\begin{enumerate}[label=(\alph*)]
  \item Soit $\gamma(t):=e^{-1/t}$ pour $t>0$ et $\gamma(t):=0$ pour $t\le0$. Montrer que $\gamma\in\Cinf(\R)$ et que $\gamma^{(k)}(0)=0$ pour tout $k$. \emph{Indication : montrer par récurrence que $\gamma^{(k)}(t)=P_k(1/t)e^{-1/t}$ pour $t>0$, avec $P_k$ un polynôme.}
  \item En déduire une fonction $\rho\in\D(\R^d)$, positive, paire, d'intégrale $1$, à support dans $\Bf(0,1)$ et strictement positive sur $B(0,1)$.
  \item (Fonction plateau en dimension 1.) Soient $a<b$ et $\varepsilon>0$. Construire explicitement $\theta\in\D(\R)$ à valeurs dans $[0,1]$, valant $1$ sur $[a,b]$ et $0$ hors de $]a-\varepsilon,b+\varepsilon[$. \emph{Indication : primitive d'une différence de deux bosses.}
  \item Montrer qu'il n'existe pas de fonction test non nulle qui soit développable en série entière au voisinage de chaque point de $\R$ (analytique). Pourquoi est-ce important pour la théorie ?
  \item Soit $\varphi\in\D(\R)$ et $\varphi_\lambda(x):=\varphi(x/\lambda)$ pour $\lambda>0$. Calculer $\norm{\varphi_\lambda}_{m,\R}$ en fonction de $\lambda$ et des normes de $\varphi$. La famille $(\varphi_\lambda)_{\lambda\to0}$ converge-t-elle dans $\D(\R)$ ? Et $(\varphi_\lambda)_{\lambda\to\infty}$ ?
\end{enumerate}
\end{exo}

\begin{solution}
\noindent\textbf{(a)} Sur $]0,\infty[$, $\gamma$ est $\Cinf$ comme composée. Montrons par récurrence que $\gamma^{(k)}(t)=P_k(1/t)e^{-1/t}$ pour $t>0$ avec $P_k$ polynomial : vrai pour $k=0$ ($P_0=1$), et si c'est vrai au rang $k$, $\gamma^{(k+1)}(t)=\big(-t^{-2}P_k'(1/t)+t^{-2}P_k(1/t)\big)e^{-1/t}=P_{k+1}(1/t)e^{-1/t}$ avec $P_{k+1}(u)=u^2(P_k(u)-P_k'(u))$. Comme $u^je^{-u}\to0$ quand $u\to+\infty$ pour tout $j$ (croissances comparées), $\gamma^{(k)}(t)\to0$ quand $t\to0^+$, pour tout $k$. On conclut par le théorème de la limite de la dérivée : $\gamma$ est continue en $0$ (car $\gamma(t)\to0=\gamma(0)$), dérivable sur $\R^\star$ avec $\gamma'(t)\to0$ des deux côtés, donc dérivable en $0$ avec $\gamma'(0)=0$ et $\gamma'$ continue ; on itère avec $\gamma'$ à la place de $\gamma$ (qui vérifie les mêmes hypothèses), d'où $\gamma\in\Cinf(\R)$ et $\gamma^{(k)}(0)=0$ pour tout $k$.

\medskip
\noindent\textbf{(b)} Posons $\rho_0(x):=\gamma(1-\norm x^2)$ pour $x\in\R^d$ : composée de $\gamma$ avec le polynôme $1-\norm x^2$, c'est une fonction $\Cinf$ ; elle est positive, paire, nulle si $\norm x\ge1$ et strictement positive si $\norm x<1$. Son support est $\Bf(0,1)$ (adhérence de la boule ouverte), compact. Son intégrale $I$ est $>0$ (fonction continue positive non identiquement nulle), et $\rho:=\rho_0/I$ convient.

\medskip
\noindent\textbf{(c)} Soit $\rho$ comme en (b) avec $d=1$, et $\rho_\varepsilon(x):=\frac2\varepsilon\rho(2x/\varepsilon)$, bosse positive d'intégrale $1$ à support dans $[-\varepsilon/2,\varepsilon/2]$ (on peut prendre n'importe quelle bosse d'intégrale $1$ à support dans $]-\varepsilon,\varepsilon[$). Posons
\[
\theta(x):=\int_{-\infty}^x\big(\rho_\varepsilon(t-a+\varepsilon/2)-\rho_\varepsilon(t-b-\varepsilon/2)\big)\,dt .
\]
La fonction intégrée est lisse à support compact, donc $\theta$ est $\Cinf$. Pour $x<a-\varepsilon$, les deux bosses sont nulles sur $]-\infty,x]$ : $\theta(x)=0$. Pour $x\in[a,b]$, la première bosse (centrée en $a-\varepsilon/2$, support dans $[a-\varepsilon,a]$) est entièrement intégrée et la seconde (support dans $[b,b+\varepsilon]$) pas encore : $\theta(x)=1-0=1$. Pour $x>b+\varepsilon$, les deux sont entièrement intégrées : $\theta(x)=1-1=0$. Ainsi $\supp\theta\subset[a-\varepsilon,b+\varepsilon]$, et $\theta\in\D(\R)$. Enfin $0\le\theta\le1$ : sur $]-\infty,a]$, $\theta$ est croissante de $0$ à $1$ (intégrale d'une fonction positive) ; sur $[b,+\infty[$, décroissante de $1$ à $0$. C'est la version « à la main » du Lemme~6.1 en dimension $1$.

\medskip
\noindent\textbf{(d)} Soit $\varphi\in\D(\R)$ analytique. Son support est compact, donc il existe $x_0$ avec $\varphi=0$ sur $]-\infty,x_0]$ (par exemple $x_0:=\min\supp\varphi$, si $\varphi\ne0$). Soit $A:=\{x\in\R:\varphi^{(k)}(x)=0\ \forall k\}$. $A$ est fermé (intersection de fermés), non vide (contient $]-\infty,x_0[$) ; il est ouvert : si $x\in A$, la série de Taylor de $\varphi$ en $x$ est nulle et, $\varphi$ étant analytique, $\varphi$ coïncide avec cette série sur un voisinage de $x$, donc y est nulle avec toutes ses dérivées. Par connexité de $\R$, $A=\R$ : $\varphi=0$.

L'importance : la théorie des distributions repose sur la possibilité de \emph{localiser} (fonctions plateaux, partitions de l'unité, Lemme~6.1 et Corollaire~6.1) et sur la densité de $\D(\Omega)$ dans tous les espaces usuels (Théorème~6.1). Les fonctions analytiques, rigides (déterminées par leur comportement en un point), ne le permettent pas ; les fonctions $\Cinf$, souples, le permettent, et (a) est exactement la preuve qu'elles sont « assez souples » : $\gamma$ est nulle sur une demi-droite sans être nulle.

\medskip
\noindent\textbf{(e)} $\varphi_\lambda^{(j)}(x)=\lambda^{-j}\varphi^{(j)}(x/\lambda)$, donc $\norm{\varphi_\lambda^{(j)}}_\infty=\lambda^{-j}\norm{\varphi^{(j)}}_\infty$ et $\norm{\varphi_\lambda}_{m,\R}=\max_{j\le m}\lambda^{-j}\norm{\varphi^{(j)}}_\infty$. Quand $\lambda\to0$ (concentration) : les supports $\lambda\supp\varphi$ restent dans un compact fixe, mais $\norm{\varphi_\lambda'}_\infty=\lambda^{-1}\norm{\varphi'}_\infty\to\infty$ dès que $\varphi$ n'est pas constante (donc, à support compact, dès que $\varphi\ne0$) : \emph{pas de convergence dans $\D(\R)$}, bien que $\varphi_\lambda\to\varphi(0)\mathbf 1_{\{0\}}$ ponctuellement. Quand $\lambda\to\infty$ (dilatation) : toutes les dérivées $\varphi_\lambda^{(j)}$, $j\ge1$, tendent uniformément vers $0$, et $\varphi_\lambda\to\varphi(0)$ uniformément sur tout compact (par continuité de $\varphi$ en $0$) ; mais les supports $\lambda\supp\varphi$ ne sont pas contenus dans un compact fixe (sauf si $\varphi=0$) : \emph{pas de convergence dans $\D(\R)$ non plus}, bien qu'il y ait convergence dans $\Cinf(\R)$ (au sens de l'Exemple~2.11) vers la constante $\varphi(0)$. Ces deux familles sont précisément celles qui servent à tester l'ordre d'une distribution (TD0 ch.~6, $\delta_0^{(k)}$) et le rôle de la condition de support dans la Définition~6.2.
\end{solution}

\begin{pointcle}
\textbf{Trois gestes sur les fonctions tests.} \emph{Translater, dilater, tronquer} : $\varphi(\cdot-a)$, $\varphi(\cdot/\lambda)$, $\theta\varphi$. Les normes $\norm\cdot_{m,K}$ se comportent de façon explicite sous ces opérations ((e) et la formule de Leibniz), et presque toutes les preuves du chapitre consistent à choisir la bonne fonction test parmi ces trois familles : une dilatation pour concentrer en un point (ordre, Dirac), une troncature pour localiser (Lemme~6.2, Exercice~6.2), une translation pour la convolution.
\end{pointcle}

% =====================================================================
\begin{exo}{Exercice 2 -- Distributions ou pas ? Ordre fini ou infini ?}
\begin{enumerate}[label=(\alph*)]
  \item Montrer que $T_1:=\sum_{k\in\Z}\delta_k$ est une distribution d'ordre $0$ sur $\R$.
  \item Montrer que $T_2:=\sum_{k\ge1}\delta_k^{(k)}$ (\ie\ $\langle T_2,\varphi\rangle=\sum_{k\ge1}(-1)^k\varphi^{(k)}(k)$) est une distribution sur $\R$, mais qu'elle n'est pas d'ordre fini.
  \item Soit $T_3(\varphi):=\sum_{k\ge1}\varphi(1/k)$. Montrer que $T_3$ définit une distribution d'ordre $0$ sur $\Omega=]0,+\infty[$, mais que la même formule ne définit pas une distribution sur $\R$ (ni même une forme linéaire sur $\D(\R)$).
  \item Montrer que $T_4(\varphi):=\sum_{k\ge1}\big(\varphi(1/k)-\varphi(0)-\tfrac1k\varphi'(0)\big)$ définit une distribution sur $\R$, d'ordre $\le2$.
\end{enumerate}
\end{exo}

\begin{solution}
\noindent\textbf{(a)} Pour $\varphi\in\D_K(\R)$, seuls les $k\in K\cap\Z$ contribuent : $N_K:=\#(K\cap\Z)<\infty$ (partie discrète d'un compact). La somme est finie, $T_1$ est linéaire, et $|\langle T_1,\varphi\rangle|\le N_K\norm\varphi_K$ : ordre $0$ (constante $C_K=N_K$, $m_K=0$).

\medskip
\noindent\textbf{(b)} \emph{Distribution :} pour $\varphi\in\D_K(\R)$, posons $N:=\max\{k\ge1:k\in K\}$ (ou $N=0$ si $K\cap\N^\star=\emptyset$). Alors $\langle T_2,\varphi\rangle=\sum_{k=1}^N(-1)^k\varphi^{(k)}(k)$ : somme finie, linéaire, et $|\langle T_2,\varphi\rangle|\le N\norm\varphi_{N,K}$. C'est (6.1) avec $m_K=N$ : $T_2$ est une distribution.

\emph{Pas d'ordre fini.} Supposons $T_2$ d'ordre $\le m$. Prenons $k:=m+1$ et $K:=[k-\frac12,k+\frac12]$. Sur $\D_K(\R)$, $\langle T_2,\varphi\rangle=(-1)^k\varphi^{(k)}(k)$ ($k$ est le seul entier de $K$, et $\varphi^{(j)}(j)=0$ pour tout entier $j\notin K$). L'hypothèse donne $C$ avec $|\varphi^{(k)}(k)|\le C\norm\varphi_{m,K}=C\norm\varphi_{k-1,K}$ pour $\varphi\in\D_K$ : c'est dire que $\delta_k^{(k)}$ est d'ordre $\le k-1$, ce qui contredit le TD0 ch.~6 (la distribution $\delta_a^{(k)}$ est d'ordre exactement $k$ ; la preuve par dilatation $\varphi_\lambda(x)=\lambda^k\varphi((x-k)/\lambda)$ est celle de l'Exercice~1(e)). Donc l'entier $m_K$ de (6.1) dépend \emph{nécessairement} de $K$ : $T_2$ est d'ordre infini.

\medskip
\noindent\textbf{(c)} \emph{Sur $\Omega=]0,\infty[$ :} un compact $K\subset]0,\infty[$ est minoré par un $\alpha>0$, donc ne contient que les $1/k$ avec $k\le1/\alpha$ : nombre fini. Pour $\varphi\in\D_K(\Omega)$, la somme est finie et $|\langle T_3,\varphi\rangle|\le\#(K\cap\{1/k\})\norm\varphi_K$ : ordre $0$. \emph{Sur $\R$ :} si $\varphi\in\D(\R)$ avec $\varphi(0)\ne0$, alors $\varphi(1/k)\to\varphi(0)\ne0$ et la série $\sum\varphi(1/k)$ diverge grossièrement : $T_3(\varphi)$ n'est pas défini. Les points $1/k$ s'accumulent en $0\in\R$, et une distribution ne peut avoir « une masse en chaque $1/k$ » sans que ces masses deviennent sommables.

\medskip
\noindent\textbf{(d)} Par la formule de Taylor avec reste intégral à l'ordre $1$ en $0$ : $\varphi(x)-\varphi(0)-x\varphi'(0)=\int_0^x(x-t)\varphi''(t)\,dt$, donc $|\varphi(1/k)-\varphi(0)-\tfrac1k\varphi'(0)|\le\frac1{2k^2}\norm{\varphi''}_\infty$. La série converge absolument, $T_4$ est bien définie et linéaire, et pour $\varphi\in\D_K(\R)$ : $|\langle T_4,\varphi\rangle|\le\frac{\pi^2}{12}\norm{\varphi''}_\infty\le\frac{\pi^2}{12}\norm\varphi_{2,K}$. C'est (6.1) avec $m_K=2$ pour tout $K$ : $T_4$ est d'ordre $\le2$. (Elle est d'ordre exactement $2$, mais nous ne le démontrons pas ; l'idée : « renormaliser » la série divergente de (c) en lui soustrayant son développement de Taylor est le procédé des \emph{parties finies}, retrouvé à l'Exercice~7.)
\end{solution}

\begin{piege}
\textbf{« Pour tout compact il existe $m_K$ » n'est pas « il existe $m$ pour tout compact ».} C'est la même inversion de quantificateurs que pour Banach--Steinhaus (TD0, Exemple~3.4). L'exemple (b) montre que la différence est réelle : $T_2$ est une distribution parfaitement légitime, mais aucun entier ne contrôle son action sur toutes les fonctions tests. Les distributions d'ordre fini sont celles qui « regardent » un nombre borné de dérivées ; une distribution à support compact est toujours d'ordre fini (TD0 ch.~6, Exemple~6.4(c)), et $T_2$ montre que l'hypothèse de support compact n'est pas superflue.
\end{piege}

% =====================================================================
\begin{exo}{Exercice 3 -- Limites au sens des distributions}
Déterminer les limites dans $\D'(\R)$, quand $n\to\infty$, des suites suivantes.
\begin{enumerate}[label=(\alph*)]
  \item $f_n:=n\mathbf 1_{[0,1/n]}$.
  \item $g_n:=n^2\big(\mathbf 1_{[0,1/n]}-\mathbf 1_{[-1/n,0]}\big)$.
  \item $S_n:=n(\delta_{1/n}-\delta_0)$.
  \item $u_n(x):=\cos^2(nx)$.
  \item $v_n(x):=\dfrac{\sin(nx)}{\pi x}$. \emph{Indication : on admettra $\int_{-\infty}^{+\infty}\frac{\sin u}u\,du=\pi$ ; séparer $\varphi(x)=\varphi(0)\chi(x)+(\varphi(x)-\varphi(0)\chi(x))$ avec $\chi$ une fonction plateau valant $1$ près de $0$, et utiliser le lemme de Riemann--Lebesgue.}
  \item Existe-t-il une suite $(\varphi_n)$ de fonctions tests convergeant vers $\delta_0$ dans $\D'(\R)$ ? convergeant vers $\delta_0$ dans $\D(\R)$ ?
\end{enumerate}
\end{exo}

\begin{solution}
Dans chaque cas on fixe $\varphi\in\D(\R)$ et l'on calcule la limite de $\langle T_n,\varphi\rangle$ ; une distribution limite est identifiée par ses valeurs sur toutes les fonctions tests.

\medskip
\noindent\textbf{(a)} $\langle f_n,\varphi\rangle=n\int_0^{1/n}\varphi=\int_0^1\varphi(t/n)\,dt\to\varphi(0)$ par convergence dominée (ou continuité uniforme). Donc $f_n\wk\delta_0$. C'est une approximation de l'unité issue de $\rho=\mathbf 1_{[0,1]}$ (Exemple~6.5), non paire : la limite est bien $\delta_0$ quand même.

\medskip
\noindent\textbf{(b)} $\langle g_n,\varphi\rangle=n^2\int_0^{1/n}\varphi-n^2\int_{-1/n}^0\varphi=n^2\int_0^{1/n}(\varphi(x)-\varphi(-x))\,dx$. Or $\varphi(x)-\varphi(-x)=2x\varphi'(0)+O(x^3)$ (Taylor à l'ordre $2$, le terme en $x^2$ étant pair ; précisément $|\varphi(x)-\varphi(-x)-2x\varphi'(0)|\le\frac{|x|^3}3\norm{\varphi'''}_\infty$). Donc $\langle g_n,\varphi\rangle=n^2\big(\varphi'(0)/n^2+O(n^{-4})\big)\to\varphi'(0)=-\langle\delta_0',\varphi\rangle$ : $g_n\wk-\delta_0'$. Plus rapide : $g_n=-F_n'$ où $F_n(x):=n(1-n|x|)^+$ est la fonction « chapeau » (continue, affine par morceaux, de pente $n^2$ sur $[-1/n,0]$ et $-n^2$ sur $[0,1/n]$ ; formule des sauts sans sauts), et $F_n$ est une approximation de l'unité (positive, d'intégrale $1$, concentrée en $0$), donc $F_n\wk\delta_0$ et, par continuité de la dérivation, $g_n=-F_n'\wk-\delta_0'$.

\medskip
\noindent\textbf{(c)} $\langle S_n,\varphi\rangle=n(\varphi(1/n)-\varphi(0))\to\varphi'(0)=-\langle\delta_0',\varphi\rangle$ : $S_n\wk-\delta_0'$. C'est cohérent avec la Proposition~6.7 appliquée à $T=\delta_0$ : comme $\tau_a\delta_0=\delta_a$ (Définition~6.7 : $\langle\tau_a\delta_0,\varphi\rangle=(\tau_{-a}\varphi)(0)=\varphi(a)$), elle donne $\frac{\tau_{-h}\delta_0-\delta_0}h=\frac{\delta_{-h}-\delta_0}h\wk\delta_0'$ quand $h\to0$ ; avec $h=-1/n$ on obtient $-n(\delta_{1/n}-\delta_0)\wk\delta_0'$, soit $S_n\wk-\delta_0'$. Attention au signe : c'est $\delta_{-h}$, et non $\delta_h$, qui apparaît dans le taux d'accroissement de la Proposition~6.7.

\medskip
\noindent\textbf{(d)} $\cos^2(nx)=\frac12+\frac12\cos(2nx)$, et $\int\cos(2nx)\varphi(x)\,dx\to0$ par Riemann--Lebesgue (TD0 ch.~6, Exemple~6.2(c)). Donc $u_n\wk\frac12$ (la fonction constante). C'est le phénomène d'\emph{homogénéisation} : une oscillation rapide converge, au sens des distributions (et faiblement dans tout $L^p_{\mathrm{loc}}$, $p<\infty$), vers sa moyenne, ici $\frac1{2\pi}\int_0^{2\pi}\cos^2=\frac12$. Noter que $u_n^2=\cos^4(nx)\wk\frac38\ne(\frac12)^2$ : la convergence faible ne passe pas aux produits.

\medskip
\noindent\textbf{(e)} Soit $\chi\in\D(\R)$ paire, valant $1$ sur $[-1,1]$, à support dans $[-2,2]$, et écrivons $\varphi=\varphi(0)\chi+\psi$ avec $\psi:=\varphi-\varphi(0)\chi\in\D(\R)$, $\psi(0)=0$. \emph{Terme en $\psi$ :} par le lemme d'Hadamard (TD0 ch.~6), $\psi(x)=x\eta(x)$ avec $\eta\in\D(\R)$, donc $\int\frac{\sin(nx)}{\pi x}\psi(x)\,dx=\frac1\pi\int\sin(nx)\eta(x)\,dx\to0$ (Riemann--Lebesgue). \emph{Terme en $\chi$ :} $\int\frac{\sin(nx)}{\pi x}\chi(x)\,dx=\frac1\pi\int_{-1}^1\frac{\sin(nx)}x\,dx+\frac1\pi\int\sin(nx)\frac{\chi(x)-\mathbf 1_{[-1,1]}(x)}x\,dx$. La fonction $\frac{\chi-\mathbf 1_{[-1,1]}}x$ est bornée (nulle sur $[-1,1]$, et $|x|\ge1$ ailleurs) à support dans $[-2,2]$, donc dans $L^1$, et le second terme tend vers $0$ par Riemann--Lebesgue dans $L^1$ (TD2 ex.~8 pour les fonctions continues, étendu à $L^1$ par densité de $\Cc$ et l'inégalité $|\int e^{inx}f|\le\norm f_1$). Le premier vaut, par le changement de variable $u=nx$, $\frac1\pi\int_{-n}^n\frac{\sin u}u\,du\to\frac1\pi\cdot\pi=1$. Au total $\langle v_n,\varphi\rangle\to\varphi(0)$ : $v_n\wk\delta_0$.

Ce noyau n'est \emph{pas} une approximation de l'unité au sens de la Définition~3.4 : $v_n$ n'est pas intégrable sur $\R$ ($|\sin(nx)/x|$ n'est pas intégrable à l'infini) et change de signe. C'est l'analogue continu du noyau de Dirichlet (TD0 ch.~6, peigne de Dirac) : sa convergence vers $\delta_0$ repose sur une compensation d'oscillations (Riemann--Lebesgue), pas sur la positivité et la concentration de la masse.

\medskip
\noindent\textbf{(f)} Oui pour $\D'$ : $\varphi_n:=\rho_n$ (Exemple~6.5), ou le Corollaire~6.2. Non pour $\D$ : si $\varphi_n\to\varphi$ dans $\D(\R)$, alors $\varphi_n\wk\varphi$ dans $\D'$ (pour $\psi\in\D$, $\int\varphi_n\psi\to\int\varphi\psi$ par convergence uniforme à supports dans un compact fixe), et par unicité de la limite dans $\D'$ (convergence simple sur les fonctions tests, dont les valeurs déterminent la distribution), $T_\varphi=\delta_0$, ce qui est exclu : $\delta_0$ n'est pas une fonction (Exemple~6.4). Plus généralement, la convergence dans $\D$ est beaucoup plus forte que la convergence dans $\D'$, et une suite de fonctions tests ne peut converger dans $\D$ que vers une fonction test.
\end{solution}

\begin{methode}
\textbf{Identifier une limite dans $\D'$.} (1) Fixer $\varphi$. (2) Mettre $\langle T_n,\varphi\rangle$ sous une forme où $n$ n'apparaît plus que dans un changement de variable ou un facteur oscillant. (3) Conclure par un des trois outils : convergence dominée (après changement de variable $x=t/n$ : concentration), développement de Taylor de $\varphi$ en un point (quand la masse se concentre et les signes se compensent : dérivées de Dirac), lemme de Riemann--Lebesgue (oscillations). (4) Reconnaître la limite : $\varphi(0)$ est $\delta_0$, $\varphi'(0)$ est $-\delta_0'$, $\int\varphi$ est la constante $1$, $\int f\varphi$ est $T_f$. Quand la suite est de la forme $T_n=S_n^{(j)}$, on peut aussi passer la dérivée à la limite (continuité de la dérivation) : c'est souvent plus rapide, cf.\ TD0 ch.~6, Exemple~6.2.
\end{methode}

\section*{Dérivation}

% =====================================================================
\begin{exo}{Exercice 4 -- Calculs de dérivées et équations différentielles}
\begin{enumerate}[label=(\alph*)]
  \item Calculer dans $\D'(\R)$ les dérivées première et seconde de $f(x)=|\sin x|$ et de $g(x)=e^{-|x|}$.
  \item En déduire que $E:=-\frac12e^{-|x|}$ vérifie $E''-E=\delta_0$ (on dit que $E$ est une \emph{solution fondamentale} de l'opérateur $\frac{d^2}{dx^2}-1$).
  \item Déterminer toutes les distributions $T\in\D'(\R)$ telles que $T''-T=\delta_0$. \emph{Indication : pour $v\in\D'(\R)$ avec $v''=v$, considérer $(e^{-x}(v'+v))'$.}
  \item Déterminer toutes les $T\in\D'(\R)$ telles que $T''+T=\delta_0$ ; vérifier que $E_1:=h(x)\sin x$ en est une.
  \item Soit $f$ continue sur $\R$, de classe $\mathcal C^1$ sur $\R\setminus\{0\}$, telle que $f'$ (définie hors de $0$) soit localement intégrable. Montrer que $T_f'=T_{f'}$. En quoi est-ce plus général que la formule des sauts ? Donner un exemple où $f'$ n'est pas bornée près de $0$.
\end{enumerate}
\end{exo}

\begin{solution}
\noindent\textbf{(a)} \emph{$f=|\sin x|$.} Elle est continue, $\mathcal C^1$ par morceaux avec points singuliers $A=\pi\Z$, et sur $]k\pi,(k+1)\pi[$, $f(x)=(-1)^k\sin x$, $f'(x)=(-1)^k\cos x=\sgn(\sin x)\cos x$. Le Théorème~6.3 est énoncé pour un ensemble $A$ \emph{fini} ; ici $A=\pi\Z$ est infini mais localement fini, et la formule des sauts reste vraie, la somme sur $A$ étant localement finie. En effet, pour $\varphi\in\D(\R)$ à support dans $[-R,R]$, soit $\theta\in\D(\R)$ un plateau valant $1$ sur $[-R-1,R+1]$ : $\theta f$ est $\mathcal C^1$ par morceaux avec un nombre fini de points singuliers (ceux de $A\cap\supp\theta$), $\langle T_f',\varphi\rangle=-\int f\varphi'=-\int\theta f\varphi'=\langle T_{\theta f}',\varphi\rangle$, et le Théorème~6.3 s'applique à $\theta f$, qui a sur $]-R-1,R+1[$ la même dérivée et les mêmes sauts que $f$. Pas de saut de $f$ : $T_f'=T_{f'}$, $f'=\sgn(\sin x)\cos x$. Cette dérivée est $\mathcal C^1$ par morceaux, de dérivée $-|\sin x|$ hors de $\pi\Z$, et de saut en $k\pi$ : $f'((k\pi)^+)-f'((k\pi)^-)=(-1)^k\cos(k\pi)-(-1)^{k-1}\cos(k\pi)=2(-1)^k(-1)^k=2$. Formule des sauts : $T_f''=-T_{|\sin x|}+2\sum_{k\in\Z}\delta_{k\pi}$. (On vérifie : $f''+f=2\sum\delta_{k\pi}$, les « cassures » de $|\sin|$ sont des pics vers le haut, d'où le signe $+$.)

\emph{$g=e^{-|x|}$.} Continue, $\mathcal C^1$ hors de $0$ avec $g'(x)=-\sgn(x)e^{-|x|}$ ; pas de saut : $T_g'=T_{g'}$. Puis $g'$ a un saut en $0$ : $g'(0^+)-g'(0^-)=-1-1=-2$, et sa dérivée hors de $0$ est $e^{-|x|}$. Donc $T_g''=T_g-2\delta_0$.

\medskip
\noindent\textbf{(b)} $E=-\frac12g$ : $E''-E=-\frac12(g''-g)=-\frac12(-2\delta_0)=\delta_0$.

\medskip
\noindent\textbf{(c)} Si $T$ et $E$ sont deux solutions, $v:=T-E$ vérifie $v''=v$. Comme $e^{-x}$ est lisse, la formule de Leibniz (Proposition~6.9) donne $(e^{-x}(v'+v))'=-e^{-x}(v'+v)+e^{-x}(v''+v')=e^{-x}(v''-v)=0$. Par le Théorème~6.4, $e^{-x}(v'+v)=c_1$ (constante), soit $v'+v=c_1e^x$ (multiplier par $e^x$, lisse, inverse de $e^{-x}$). De même $(e^xv)'=e^x(v'+v)=c_1e^{2x}$, et $e^xv-\frac{c_1}2e^{2x}$ a une dérivée nulle : $e^xv=\frac{c_1}2e^{2x}+c_2$, soit $v=\frac{c_1}2e^x+c_2e^{-x}$. Réciproquement ces fonctions lisses vérifient $v''=v$. Les solutions de $T''-T=\delta_0$ sont donc
\[
T=-\tfrac12e^{-|x|}+ae^x+be^{-x},\qquad a,b\in\K .
\]
\emph{Toutes les solutions distributions de l'équation homogène sont des fonctions lisses} : c'est un cas très simple de régularité elliptique (comparer avec le lemme de Weyl, Lemme~6.4).

\medskip
\noindent\textbf{(d)} $E_1=h\sin x$ est continue (car $\sin0=0$), $\mathcal C^1$ hors de $0$ avec $E_1'=h\cos x$ (pas de saut de $E_1$), puis $E_1'$ a un saut $\cos0-0=1$ en $0$ et pour dérivée $-h\sin x$ hors de $0$ : $E_1''=-h\sin x+\delta_0=-E_1+\delta_0$. Donc $E_1''+E_1=\delta_0$. Pour l'équation homogène $v''+v=0$ : on utilise la même astuce avec $e^{\pm ix}$ (on se place dans $\K=\C$ ; pour $\K=\R$, raisonner sur les parties réelle et imaginaire, ou utiliser $(\cos x\,v'+\sin x\,v)'=\cos x\,(v''+v)=0$ et $(\sin x\,v'-\cos x\,v)'=\sin x\,(v''+v)=0$, puis résoudre le système linéaire $\cos x\,v'+\sin x\,v=c_1$, $\sin x\,v'-\cos x\,v=c_2$, de déterminant $-1$ : $v=c_1\sin x-c_2\cos x$). Les solutions sont $T=h(x)\sin x+a\cos x+b\sin x$.

\medskip
\noindent\textbf{(e)} Pour $\varphi\in\D(\R)$ et $\varepsilon>0$ : $\int_\R f\varphi'=\lim_{\varepsilon\to0}\big(\int_{-\infty}^{-\varepsilon}f\varphi'+\int_\varepsilon^\infty f\varphi'\big)$ (convergence dominée, $f\varphi'\in L^1$). Sur $[\varepsilon,R]$ (où $\supp\varphi\subset[-R,R]$), $f$ est $\mathcal C^1$ : $\int_\varepsilon^Rf\varphi'=-f(\varepsilon)\varphi(\varepsilon)-\int_\varepsilon^Rf'\varphi$, et de même $\int_{-R}^{-\varepsilon}f\varphi'=f(-\varepsilon)\varphi(-\varepsilon)-\int_{-R}^{-\varepsilon}f'\varphi$. Les termes de bord tendent vers $f(0)\varphi(0)-f(0)\varphi(0)=0$ par continuité de $f$ (c'est là que l'hypothèse « $f$ continue en $0$ » sert : pas de saut), et $\int_{\varepsilon<|x|<R}f'\varphi\to\int f'\varphi$ par convergence dominée puisque $f'\varphi\in L^1$. Donc $\langle T_f',\varphi\rangle=-\int f\varphi'=\int f'\varphi=\langle T_{f'},\varphi\rangle$.

C'est plus général que le Théorème~6.3, qui exige des limites finies de $f'$ à droite et à gauche (donc $f'$ localement bornée). Exemple : $f(x)=\sqrt{|x|}$, continue, $f'(x)=\frac{\sgn x}{2\sqrt{|x|}}$ hors de $0$, non bornée mais localement intégrable (intégrale de $|x|^{-1/2}$ convergente) : $T_f'=T_{f'}$. Autre exemple : $f(x)=x^2\sin(1/x)$ ($f(0)=0$), dont la dérivée $f'(x)=2x\sin(1/x)-\cos(1/x)$ est bornée mais n'a de limite ni à droite ni à gauche en $0$. (En revanche $x\sin(1/x)$ ne vérifie pas l'hypothèse : sa dérivée $\sin(1/x)-\frac1x\cos(1/x)$ n'est pas localement intégrable, car $\int_0^1\frac{|\cos(1/x)|}x\,dx=\int_1^{+\infty}\frac{|\cos u|}u\,du=+\infty$.) Et un contre-exemple instructif : $f=\ln|x|$ n'est pas continue en $0$ mais la conclusion reste vraie « à condition de lire $f'=1/x$ comme $\mathrm{vp}(1/x)$ » (Proposition~6.4) — le terme de bord $\ln\varepsilon\,(\varphi(-\varepsilon)-\varphi(\varepsilon))\to0$ survit à la discontinuité, parce qu'elle est symétrique.
\end{solution}

\begin{pointcle}
\textbf{Solution fondamentale.} Une solution fondamentale d'un opérateur différentiel $P(\frac d{dx})$ est une distribution $E$ avec $P(\frac d{dx})E=\delta_0$. On en construit en dimension $1$ par la formule des sauts : $E=h\cdot y$ où $y$ est la solution de l'équation homogène avec $y(0)=\dots=y^{(n-2)}(0)=0$, $y^{(n-1)}(0)=1$ (ici $y=\sinh x$ pour $\frac{d^2}{dx^2}-1$, soit $E=h\sinh x$, qui diffère de $-\frac12e^{-|x|}$ par la solution homogène $\frac12e^x$ ; et $y=\sin x$ pour $\frac{d^2}{dx^2}+1$). L'intérêt est à l'Exercice~8 : $u=E\star\psi$ résout $P(\frac d{dx})u=\psi$.
\end{pointcle}

% =====================================================================
\begin{exo}{Exercice 5 -- Primitives d'une distribution}
Soit $I\subset\R$ un intervalle ouvert et $T\in\D'(I)$.
\begin{enumerate}[label=(\alph*)]
  \item Soit $\theta\in\D(I)$ d'intégrale $1$. Montrer que pour toute $\varphi\in\D(I)$, la fonction $\psi:=\varphi-\big(\int_I\varphi\big)\theta$ admet une unique primitive $\Psi\in\D(I)$, et que $\varphi\mapsto\Psi$ est linéaire et vérifie $\norm\Psi_{m+1,K'}\le C_K\norm\varphi_{m,K'}$ pour $\varphi\in\D_K(I)$, avec $K':=\mathrm{conv}(K\cup\supp\theta)$ (enveloppe convexe : un segment de $I$) et des constantes convenables (préciser).
  \item On pose $\langle S,\varphi\rangle:=-\langle T,\Psi\rangle$. Montrer que $S\in\D'(I)$ et que $S'=T$.
  \item Montrer que les primitives de $T$ (les $U\in\D'(I)$ avec $U'=T$) sont exactement les $S+c$, $c\in\K$.
  \item Application : soit $g\in L^1_{\mathrm{loc}}(I)$ et $U\in\D'(I)$ avec $U'=T_g$. Montrer que $U$ est une fonction continue, de la forme $U=T_F$ avec $F(x)=c+\int_{x_0}^xg$.
  \item Montrer qu'une primitive d'une distribution d'ordre $m\ge1$ est d'ordre $\le m-1$, et qu'une primitive d'une distribution d'ordre $0$ est d'ordre $0$. Quel est l'ordre des primitives de $\delta_0$ ? de $\delta_0'$ ?
\end{enumerate}
\end{exo}

\begin{solution}
\noindent\textbf{(a)} $\psi\in\D(I)$ et $\int_I\psi=\int\varphi-\int\varphi=0$. Posons $\Psi(x):=\int_{-\infty}^x\psi(t)\,dt$ (intégrale sur $I\cap]-\infty,x]$, prolongée par $0$ hors de $I$) : $\Psi$ est $\Cinf$ avec $\Psi'=\psi$, nulle à gauche du support de $\psi$ ; et à droite de ce support, $\Psi(x)=\int_I\psi=0$. Donc $\supp\Psi\subset\mathrm{conv}(\supp\psi)\subset K'$, compact de $I$ (l'enveloppe convexe d'un compact de l'intervalle $I$ est un segment contenu dans $I$) : $\Psi\in\D(I)$. Unicité : deux primitives diffèrent d'une constante, nulle puisque à support compact. Linéarité : $\varphi\mapsto\psi\mapsto\Psi$ est composée de deux applications linéaires. Estimation : $\Psi^{(j+1)}=\psi^{(j)}=\varphi^{(j)}-(\int\varphi)\theta^{(j)}$, donc $\norm{\Psi^{(j+1)}}_\infty\le\norm{\varphi^{(j)}}_\infty+|K|\,\norm\varphi_\infty\norm{\theta^{(j)}}_\infty$ ($|K|$ désignant la longueur de $\mathrm{conv}K$) ; et $\norm\Psi_\infty\le|K'|\,\norm\psi_\infty\le|K'|(1+|K|\norm\theta_\infty)\norm\varphi_\infty$. D'où $\norm\Psi_{m+1,K'}\le C_{K,\theta,m}\norm\varphi_{m,K'}$, et même $\norm\Psi_{m+1,K'}\le C\norm\varphi_{m,K}$ puisque $\varphi$ est à support dans $K$.

\medskip
\noindent\textbf{(b)} $S$ est linéaire (composée). Continuité : pour $\varphi\in\D_K(I)$, $\Psi\in\D_{\tilde K}(I)$ avec $\tilde K:=\mathrm{conv}(K\cup\supp\theta)$, et (6.1) pour $T$ sur $\tilde K$ donne $|\langle S,\varphi\rangle|\le C_{\tilde K}\norm\Psi_{m_{\tilde K},\tilde K}\le C_{\tilde K}C'\norm\varphi_{m_{\tilde K},K}$ (avec (a), en notant qu'une majoration de $\norm\Psi_{m+1}$ en majore aussi $\norm\Psi_m$). Donc $S\in\D'(I)$. Calculons $S'$ : $\langle S',\varphi\rangle=-\langle S,\varphi'\rangle=\langle T,\Psi_1\rangle$ où $\Psi_1$ est la primitive à support compact de $\varphi'-(\int\varphi')\theta=\varphi'$ (car $\int\varphi'=0$) : c'est $\Psi_1=\varphi$. Donc $\langle S',\varphi\rangle=\langle T,\varphi\rangle$ : $S'=T$.

\medskip
\noindent\textbf{(c)} Si $U'=T=S'$, alors $(U-S)'=0$ sur l'intervalle (connexe) $I$, donc $U-S=c$ par le Théorème~6.4. Réciproquement $(S+c)'=S'$.

\medskip
\noindent\textbf{(d)} Soit $x_0\in I$ et $G(x):=\int_{x_0}^xg$ : fonction continue (convergence dominée). Par (b) et (c), $U=S+c$ où $S$ est la primitive construite en (b) à partir de $T=T_g$ : il suffit donc de \emph{calculer $S$}. Soit $\varphi\in\D(I)$, $\psi:=\varphi-(\int\varphi)\theta$ et $\Psi$ sa primitive à support compact ; soit $[a,b]\subset I$ un segment contenant $\supp\psi$ (donc $\supp\Psi$), de sorte que $\Psi(x)=\int_a^x\psi(t)\,dt$ pour $x\in[a,b]$. Alors
\[
\langle S,\varphi\rangle=-\langle T_g,\Psi\rangle=-\int_a^bg(x)\Big(\int_a^x\psi(t)\,dt\Big)dx=-\int_a^b\psi(t)\Big(\int_t^bg(x)\,dx\Big)dt ,
\]
l'échange étant licite par Fubini sur $[a,b]^2$ ($g\in L^1([a,b])$ et $\psi$ bornée). Or $\int_t^bg=G(b)-G(t)$ et $\int\psi=0$, donc le terme en $G(b)$ disparaît :
\[
\langle S,\varphi\rangle=\int_I\psi\,G=\int_I\varphi\,G-\Big(\int_I\varphi\Big)\Big(\int_I\theta G\Big)=\langle T_{G-c_\theta},\varphi\rangle,\qquad c_\theta:=\int_I\theta G .
\]
Ainsi $S=T_{G-c_\theta}$ est une fonction continue, et $U=T_{G+c'}$ avec $c'=c-c_\theta$. Au passage, $S'=T_g$ donne $T_G'=T_g$ : c'est le théorème fondamental de l'analyse au sens des distributions, pour $g$ seulement localement intégrable, et l'on voit que le choix de $\theta$ ne fait que fixer la constante d'intégration (par la normalisation $\int\theta\,(G-c_\theta)=0$). En particulier \emph{une distribution dont la dérivée est localement intégrable est une fonction continue} — et même absolument continue.

\emph{Variante par régularisation}, qui établit directement $T_G'=T_g$ sans passer par (b) : fixons $\varphi\in\D(I)$ et un segment $[a,b]\subset I$ contenant $x_0$ et $\supp\varphi$ ; soit $g_n\in\Cinf(\R)$ avec $g_n\to g$ dans $L^1([a,b])$ (régulariser $g\mathbf 1_{[a,b]}$ par $\rho_n$, Théorème~3.6) et $G_n(x):=\int_{x_0}^xg_n$, de classe $\mathcal C^1$ avec $G_n'=g_n$ et $|G_n(x)-G(x)|\le\norm{g_n-g}_{L^1([a,b])}\to0$ uniformément sur $[a,b]$. L'intégration par parties classique donne $-\int G_n\varphi'=\int g_n\varphi$, et l'on passe à la limite des deux côtés : $-\int G\varphi'=\int g\varphi$, soit $T_G'=T_g$ ; puis $U=T_G+c$ par (c).

\medskip
\noindent\textbf{(e)} L'estimation de (b) donne, si $T$ est d'ordre $m$ : $|\langle S,\varphi\rangle|\le C\norm\Psi_{m,\tilde K}$ ; or $\norm\Psi_{m,\tilde K}$ ne fait intervenir que $\Psi,\dots,\Psi^{(m)}$, c'est-à-dire $\Psi$ et $\psi,\dots,\psi^{(m-1)}$, contrôlées par $\norm\varphi_{m-1,K}$ si $m\ge1$ (et par $\norm\varphi_\infty$ si $m=0$). Donc \emph{une primitive d'une distribution d'ordre $m\ge1$ est d'ordre $\le m-1$}, et une primitive d'une distribution d'ordre $0$ est d'ordre $0$ ; l'intégration \emph{régularise}, d'un cran. Exemples : les primitives de $\delta_0$ sont $h+c$, d'ordre $0$ (fonctions) ; celles de $\delta_0'$ sont $\delta_0+c$, d'ordre $0$. En sens inverse, la dérivation fait monter l'ordre d'au plus $1$ (Proposition~6.3), et exactement de $1$ pour $\delta_0^{(k)}$ (TD0 ch.~6).
\end{solution}

\begin{commentaire}
C'est le Lemme~6.3 en dimension $1$ (les fonctions tests d'intégrale nulle sont exactement les dérivées de fonctions tests), transformé en \emph{construction} d'une primitive : on ne peut pas poser « $\langle S,\varphi\rangle=-\langle T,\int\varphi\rangle$ » car $\int_{-\infty}^x\varphi$ n'est pas à support compact ; on corrige $\varphi$ par un multiple de $\theta$ pour rendre l'intégrale nulle, et c'est le Théorème~6.4 qui garantit que le choix de $\theta$ ne change $S$ que d'une constante. Le point (d) est la version distributionnelle du théorème fondamental de l'analyse : il dit que l'équation $U'=g$ n'a pas d'autres solutions dans $\D'$ que les solutions « classiques » (absolument continues), même pour $g$ seulement localement intégrable — à comparer avec l'Exemple~6.10 et le TD0 ch.~6.
\end{commentaire}

\section*{Multiplication, support, valeur principale}

% =====================================================================
\begin{exo}{Exercice 6 -- Équations $xT=\dots$ et distributions à support ponctuel}
On rappelle (TD0 ch.~6, lemme d'Hadamard) que si $\varphi\in\D(\R)$ vérifie $\varphi(0)=0$, alors $\varphi=x\psi$ avec $\psi\in\D(\R)$, et que les solutions de $xT=0$ sont les $c\delta_0$.
\begin{enumerate}[label=(\alph*)]
  \item Résoudre dans $\D'(\R)$ : $xT=\delta_0$ ; puis $x^2T=0$ ; puis $x^kT=0$ ($k\ge1$).
  \item Soient $a\ne b$. Résoudre $(x-a)(x-b)T=0$. \emph{Indication : écrire $T=\theta T+(1-\theta)T$ avec $\theta\in\D(\R)$ valant $1$ près de $a$ et $0$ près de $b$, et observer que $x-b$ est « inversible » sur le support de $\theta$.}
  \item On dit que $T\in\D'(\Omega)$ est \emph{nulle sur l'ouvert} $U\subset\Omega$ si $\langle T,\varphi\rangle=0$ pour toute $\varphi\in\D(U)$. Montrer, à l'aide du Corollaire~6.1, que si $T$ est nulle sur chaque ouvert d'une famille $(U_i)_{i\in I}$, elle est nulle sur $\bigcup_iU_i$. On définit le \emph{support} de $T$ comme le complémentaire du plus grand ouvert sur lequel $T$ est nulle. Déterminer les supports de $T_f$ ($f$ continue), de $\delta_a^{(k)}$, de $\mathrm{vp}(1/x)$.
  \item Soit $T\in\D'(\R)$ de support $\{0\}$ et d'ordre $\le m$ (une distribution à support compact est d'ordre fini, TD0 ch.~6). Montrer que $x^{m+1}T=0$, puis que $T=\sum_{j=0}^mc_j\delta_0^{(j)}$. \emph{Indication : pour $\varphi\in\D(\R)$, écrire $\langle T,x^{m+1}\varphi\rangle=\langle T,\chi_\varepsilon x^{m+1}\varphi\rangle$ avec $\chi_\varepsilon(x)=\chi(x/\varepsilon)$, et faire $\varepsilon\to0$.}
\end{enumerate}
\end{exo}

\begin{solution}
\noindent\textbf{(a)} \emph{$xT=\delta_0$.} Une solution particulière : $x\delta_0'=-\delta_0$ (TD0 ch.~6), donc $T_0:=-\delta_0'$ convient. Les solutions sont $T_0+$ (solutions de $xT=0$) $=-\delta_0'+c\delta_0$.

\emph{$x^2T=0$.} Posons $U:=xT$ ; alors $xU=0$, donc $U=c\delta_0$, soit $xT=c\delta_0$ : par ce qui précède, $T=-c\delta_0'+c'\delta_0$. Réciproquement $x^2\delta_0=0$ et $x^2\delta_0'=x(x\delta_0')=-x\delta_0=0$. Les solutions sont $\Vect(\delta_0,\delta_0')$.

\emph{$x^kT=0$.} Par récurrence, les solutions sont $\Vect(\delta_0,\dots,\delta_0^{(k-1)})$ : si $x^kT=0$, alors $U:=xT$ vérifie $x^{k-1}U=0$, donc $U=\sum_{j<k-1}a_j\delta_0^{(j)}$ ; comme $x\delta_0^{(j+1)}=-(j+1)\delta_0^{(j)}$ (Leibniz : $(x\delta_0^{(j)})'=\delta_0^{(j)}+x\delta_0^{(j+1)}$, et par récurrence $x\delta_0^{(j)}=-j\delta_0^{(j-1)}$), une solution particulière de $xT=U$ est $-\sum_j\frac{a_j}{j+1}\delta_0^{(j+1)}$, et l'on ajoute $c\delta_0$. Réciproquement $x^k\delta_0^{(j)}=0$ pour $j<k$ : $\langle x^k\delta_0^{(j)},\varphi\rangle=(-1)^j(x^k\varphi)^{(j)}(0)=0$ car toutes les dérivées d'ordre $<k$ de $x^k\varphi$ s'annulent en $0$ (Leibniz).

\medskip
\noindent\textbf{(b)} Soit $\theta\in\D(\R)$ valant $1$ sur un voisinage $]a-r,a+r[$ de $a$ et à support dans $]a-2r,a+2r[$, avec $2r<|b-a|$ (Lemme~6.1 ou Exercice~1(c)). Écrivons $T=\theta T+(1-\theta)T$.

\emph{Le morceau $\theta T$.} On a $(x-a)(x-b)\theta T=\theta(x-a)(x-b)T=0$ (la multiplication par des fonctions lisses est commutative et associative : $\langle fgT,\varphi\rangle=\langle T,fg\varphi\rangle$). Soit $\eta\in\Cinf(\R)$ égale à $\frac1{x-b}$ sur $[a-2r,a+2r]$ (prolonger arbitrairement de façon lisse hors de cet intervalle, par exemple $\eta:=\tilde\theta/(x-b)$ avec $\tilde\theta$ plateau valant $1$ sur $[a-2r,a+2r]$ et nulle près de $b$). Alors $\eta(x-b)\theta=\theta$ (égalité de fonctions : vraie sur $\supp\theta$ où $\eta(x-b)=1$, et les deux membres sont nuls ailleurs), donc $(x-a)\theta T=\eta(x-b)(x-a)\theta T=0$. Par translation ($\tau_{-a}$ transforme $(x-a)S=0$ en $xS=0$, Exercice~6.3), $\theta T=\alpha\delta_a$.

\emph{Le morceau $(1-\theta)T$.} De même, $(1-\theta)$ est nulle sur $]a-r,a+r[$, et $x-a$ est inversible sur le fermé $\{|x-a|\ge r\}\supset\supp(1-\theta)$ : avec $\zeta\in\Cinf$ égale à $\frac1{x-a}$ sur $\{|x-a|\ge r\}$, on a $\zeta(x-a)(1-\theta)=(1-\theta)$ et $(x-b)(1-\theta)T=\zeta(x-a)(x-b)(1-\theta)T=0$, donc $(1-\theta)T=\beta\delta_b$.

Conclusion : $T=\alpha\delta_a+\beta\delta_b$, et réciproquement ces distributions conviennent ($(x-a)\delta_a=0$, TD0 ch.~6). Les solutions forment un espace de dimension $2$, comme pour l'équation $x^2T=0$ mais avec deux Dirac distincts au lieu de $\delta_0,\delta_0'$ : quand $b\to a$, $\frac{\delta_b-\delta_a}{b-a}\wk-\delta_a'$ (Exercice~3(c)), et $\Vect(\delta_a,\delta_b)$ « tend » vers $\Vect(\delta_a,\delta_a')$ : c'est cohérent.

\medskip
\noindent\textbf{(c)} Soit $\varphi\in\D(\bigcup_iU_i)$, de support compact $K$. Par compacité, $K\subset U_{i_1}\cup\dots\cup U_{i_N}$. Le Corollaire~6.1 fournit $\eta_k\in\D(U_{i_k})$, $1\le k\le N$, avec $\sum\eta_k=1$ sur un voisinage de $K$. Alors $\varphi=\sum_k\eta_k\varphi$ (égalité sur le voisinage de $K$, et les deux membres sont nuls hors de $K$), avec $\eta_k\varphi\in\D(U_{i_k})$, donc $\langle T,\varphi\rangle=\sum_k\langle T,\eta_k\varphi\rangle=0$. Ainsi la réunion de tous les ouverts où $T$ est nulle est encore un ouvert où $T$ est nulle : c'est le plus grand, et le support est bien défini (fermé de $\Omega$).

\emph{Supports.} $\supp T_f=\supp f=\overline{\{f\ne0\}}$ pour $f$ continue : si $U$ est un ouvert où $f\ne0$ quelque part, prendre $\varphi\ge0$ une bosse autour d'un point où $f\ne0$ (assez petite pour que $f$ y garde un signe constant, par continuité) : $\int f\varphi\ne0$ ; et $T_f$ est nulle sur l'ouvert $\{f=0\}^\circ$ évidemment. $\supp\delta_a^{(k)}=\{a\}$ : nulle sur $\R\setminus\{a\}$, pas sur un ouvert contenant $a$ (prendre $\varphi$ avec $\varphi^{(k)}(a)\ne0$). $\supp\mathrm{vp}(1/x)=\R$ : sur tout ouvert $U$, si $U\not\ni0$, $\mathrm{vp}(1/x)$ y coïncide avec la fonction $1/x$ qui ne s'annule pas ; si $0\in U$, $U$ contient un intervalle $]0,r[$ où l'argument précédent s'applique.

\emph{Support compact et ordre fini.} Si $\supp T$ est compact, $T$ est d'ordre fini (fait utilisé en (d)). Soit $\chi\in\D(\Omega)$ valant $1$ sur un voisinage de $\supp T$ (Lemme~6.1). Pour $\varphi\in\D(\Omega)$, $(1-\chi)\varphi$ est une fonction test à support dans l'ouvert $\Omega\setminus\supp T$, sur lequel $T$ est nulle : $\langle T,\varphi\rangle=\langle T,\chi\varphi\rangle$. L'estimation (6.1) sur le compact $L:=\supp\chi$ et la formule de Leibniz donnent alors $|\langle T,\varphi\rangle|\le C_L\norm{\chi\varphi}_{m_L,L}\le C\norm\varphi_{m_L,L}$ pour \emph{toute} $\varphi\in\D(\Omega)$ : $T$ est d'ordre $\le m_L$, et c'est exactement l'hypothèse du TD0 ch.~6, Exemple~6.4(c).

\medskip
\noindent\textbf{(d)} \emph{$x^{m+1}T=0$.} Soit $\chi\in\D(]-1,1[)$ valant $1$ sur $[-\frac12,\frac12]$, $\chi_\varepsilon:=\chi(\cdot/\varepsilon)$, et $\varphi\in\D(\R)$. La fonction $(1-\chi_\varepsilon)x^{m+1}\varphi$ est une fonction test nulle sur $]-\varepsilon/2,\varepsilon/2[$, donc à support dans $\R\setminus\{0\}$ où $T$ est nulle : $\langle T,x^{m+1}\varphi\rangle=\langle T,\chi_\varepsilon x^{m+1}\varphi\rangle$ pour tout $\varepsilon\in]0,1]$. Comme $T$ est d'ordre $\le m$ et que $\chi_\varepsilon x^{m+1}\varphi\in\D_{[-1,1]}$, il existe $C$ (indépendante de $\varepsilon$) avec $|\langle T,x^{m+1}\varphi\rangle|\le C\norm{\chi_\varepsilon x^{m+1}\varphi}_{m,[-1,1]}$. Estimons, pour $j\le m$, par Leibniz :
\[
(\chi_\varepsilon x^{m+1}\varphi)^{(j)}=\sum_{i=0}^j\binom ji\varepsilon^{-i}\chi^{(i)}(x/\varepsilon)\,(x^{m+1}\varphi)^{(j-i)}(x).
\]
Sur le support de $\chi_\varepsilon$, $|x|\le\varepsilon$, et $(x^{m+1}\varphi)^{(j-i)}$ est une combinaison de termes $x^{m+1-l}\varphi^{(j-i-l)}$ ($l\le j-i\le m$), chacun majoré par $\varepsilon^{m+1-l}\norm\varphi_{m,\R}\cdot(m+1)!$ ; donc $|(x^{m+1}\varphi)^{(j-i)}(x)|\le C'\varepsilon^{m+1-(j-i)}$ pour $|x|\le\varepsilon$. Au total $|(\chi_\varepsilon x^{m+1}\varphi)^{(j)}|\le C''\sum_i\varepsilon^{-i}\varepsilon^{m+1-j+i}=C'''\varepsilon^{m+1-j}\le C'''\varepsilon$ pour $j\le m$. Donc $|\langle T,x^{m+1}\varphi\rangle|\le C\,C'''\varepsilon\to0$ : $x^{m+1}T=0$.

\emph{Conclusion.} Par (a), $T\in\Vect(\delta_0,\dots,\delta_0^{(m)})$. (On peut aussi conclure directement : par Taylor et Hadamard itéré, $\varphi=\sum_{j\le m}\frac{\varphi^{(j)}(0)}{j!}x^j\chi+x^{m+1}\psi$ avec $\psi\in\D$, d'où $\langle T,\varphi\rangle=\sum_j\frac{\langle T,x^j\chi\rangle}{j!}\varphi^{(j)}(0)$.)
\end{solution}

\begin{pointcle}
\textbf{Localiser avec une fonction plateau.} Pour résoudre une équation $fT=0$ avec $f$ lisse ayant des zéros isolés, on découpe $T$ par une partition de l'unité adaptée aux zéros ; sur chaque morceau, $f$ est le produit d'un facteur « inversible » (que l'on absorbe grâce à une fonction lisse $\eta$ qui coïncide avec $1/f$ sur le support du morceau) et du facteur $(x-a)^k$, et l'on est ramené à $x^kS=0$. Deux faits rendent cela licite : la multiplication par les fonctions lisses est commutative et associative, et une fonction lisse n'agit sur $\theta T$ qu'à travers ses valeurs sur $\supp\theta$. Le résultat (d) — \emph{une distribution à support ponctuel est une combinaison finie de dérivées de Dirac} — est le théorème de structure le plus utile du chapitre ; l'argument de dilatation $\chi_\varepsilon$ est le même que pour l'ordre de $\delta_0^{(k)}$, lu à l'envers.
\end{pointcle}

% =====================================================================
\begin{exo}{Exercice 7 -- Valeur principale, partie finie, formule de Sokhotski--Plemelj}
\begin{enumerate}[label=(\alph*)]
  \item On pose $\mathrm{Pf}(1/x^2):=-\big(\mathrm{vp}(1/x)\big)'$. Montrer que pour $\varphi\in\D(\R)$,
  \[
  \big\langle\mathrm{Pf}(1/x^2),\varphi\big\rangle=\lim_{\varepsilon\to0}\Big(\int_{|x|>\varepsilon}\frac{\varphi(x)}{x^2}\,dx-\frac{2\varphi(0)}\varepsilon\Big)=\int_\R\frac{\varphi(x)-\varphi(0)-x\varphi'(0)\mathbf 1_{|x|<1}}{x^2}\,dx .
  \]
  \item Montrer que $x\,\mathrm{Pf}(1/x^2)=\mathrm{vp}(1/x)$ et $x^2\,\mathrm{Pf}(1/x^2)=1$. Quelles sont toutes les solutions de $x^2T=1$ ?
  \item Résoudre dans $\D'(\R)$ l'équation $xT'+T=0$.
  \item (Sokhotski--Plemelj.) Pour $\varepsilon>0$, soit $f_\varepsilon(x):=\dfrac1{x-i\varepsilon}$. Montrer que $f_\varepsilon\in L^1_{\mathrm{loc}}(\R)$ et que
  \[
  f_\varepsilon\xrightarrow[\varepsilon\to0^+]{}\mathrm{vp}(1/x)+i\pi\delta_0\quad\text{dans }\D'(\R).
  \]
  En déduire la limite de $\frac1{x+i\varepsilon}$, et retrouver $x\,\mathrm{vp}(1/x)=1$.
  \item Montrer que $\mathrm{vp}(1/x)$ est la limite dans $\D'(\R)$ de $\dfrac{x}{x^2+\varepsilon^2}$, et que $\delta_0$ est la limite de $\dfrac1\pi\dfrac{\varepsilon}{x^2+\varepsilon^2}$ (noyau de Poisson).
\end{enumerate}
\end{exo}

\begin{solution}
\noindent\textbf{(a)} Par définition, $\langle\mathrm{Pf}(1/x^2),\varphi\rangle=\langle\mathrm{vp}(1/x),\varphi'\rangle=\lim_{\varepsilon\to0}\int_{|x|>\varepsilon}\frac{\varphi'(x)}x\,dx$. Intégrons par parties sur $[\varepsilon,R]$ et $[-R,-\varepsilon]$ ($\supp\varphi\subset[-R,R]$) :
\[
\int_\varepsilon^R\frac{\varphi'(x)}x\,dx=\Big[\frac{\varphi(x)}x\Big]_\varepsilon^R+\int_\varepsilon^R\frac{\varphi(x)}{x^2}\,dx=-\frac{\varphi(\varepsilon)}\varepsilon+\int_\varepsilon^R\frac{\varphi}{x^2},\qquad
\int_{-R}^{-\varepsilon}\frac{\varphi'(x)}x\,dx=\frac{\varphi(-\varepsilon)}{-\varepsilon}+\int_{-R}^{-\varepsilon}\frac\varphi{x^2}.
\]
Somme : $\int_{|x|>\varepsilon}\frac\varphi{x^2}-\frac{\varphi(\varepsilon)+\varphi(-\varepsilon)}\varepsilon$. Or $\varphi(\varepsilon)+\varphi(-\varepsilon)=2\varphi(0)+O(\varepsilon^2)$ (Taylor à l'ordre $2$, les termes impairs se compensant), donc $\frac{\varphi(\varepsilon)+\varphi(-\varepsilon)}\varepsilon=\frac{2\varphi(0)}\varepsilon+O(\varepsilon)$, et la première formule suit. Pour la seconde : $\int_{|x|>\varepsilon}\frac{\varphi(x)-\varphi(0)-x\varphi'(0)\mathbf 1_{|x|<1}}{x^2}\,dx=\int_{|x|>\varepsilon}\frac\varphi{x^2}-\varphi(0)\int_{|x|>\varepsilon}\frac{dx}{x^2}-\varphi'(0)\int_{\varepsilon<|x|<1}\frac{dx}x$ ; la dernière intégrale est nulle (imparité), $\int_{|x|>\varepsilon}x^{-2}\,dx=2/\varepsilon$, et l'intégrande de gauche est intégrable sur $\R$ (il est $O(1)$ près de $0$ par Taylor, et $O(x^{-2})$ à l'infini) : on peut faire $\varepsilon\to0$ par convergence dominée, ce qui donne l'égalité des deux expressions.

\medskip
\noindent\textbf{(b)} $\langle x\,\mathrm{Pf}(1/x^2),\varphi\rangle=\langle\mathrm{Pf}(1/x^2),x\varphi\rangle=\lim_\varepsilon\big(\int_{|x|>\varepsilon}\frac{x\varphi}{x^2}-\frac{2\cdot0}\varepsilon\big)=\lim_\varepsilon\int_{|x|>\varepsilon}\frac\varphi x=\langle\mathrm{vp}(1/x),\varphi\rangle$ (la fonction $x\varphi$ s'annule en $0$). Puis $x^2\mathrm{Pf}(1/x^2)=x\,\mathrm{vp}(1/x)=1$. Les solutions de $x^2T=1$ sont $\mathrm{Pf}(1/x^2)+a\delta_0+b\delta_0'$ (Exercice~6(a)). Autre vérification de $x\,\mathrm{Pf}(1/x^2)=\mathrm{vp}(1/x)$ par Leibniz : $(x\,\mathrm{vp}(1/x))'=\mathrm{vp}(1/x)+x\,\mathrm{vp}(1/x)'$, et le membre de gauche est $1'=0$, donc $x\,\mathrm{vp}(1/x)'=-\mathrm{vp}(1/x)$, soit $x\,\mathrm{Pf}(1/x^2)=\mathrm{vp}(1/x)$.

\medskip
\noindent\textbf{(c)} Par Leibniz, $xT'+T=(xT)'$. L'équation équivaut à $(xT)'=0$, soit $xT=c$ (Théorème~6.4), soit $T=c\,\mathrm{vp}(1/x)+d\,\delta_0$ (TD0 ch.~6 : solutions de $xT=1$ et de $xT=0$). Vérification directe : $x\,\delta_0'+\delta_0=0$ et $x\,\mathrm{vp}(1/x)'+\mathrm{vp}(1/x)=-\mathrm{vp}(1/x)+\mathrm{vp}(1/x)=0$. C'est l'équation d'Euler $xy'+y=0$, dont les solutions classiques sur $\R^\star$ sont $c/x$ ; au sens des distributions, $1/x$ doit être lu comme $\mathrm{vp}(1/x)$, et il apparaît une solution supplémentaire $\delta_0$, « concentrée » au point singulier de l'équation.

\medskip
\noindent\textbf{(d)} $|f_\varepsilon(x)|=(x^2+\varepsilon^2)^{-1/2}\le1/\varepsilon$ : $f_\varepsilon$ est continue bornée, donc $L^1_{\mathrm{loc}}$. Écrivons $\frac1{x-i\varepsilon}=\frac{x+i\varepsilon}{x^2+\varepsilon^2}=\frac x{x^2+\varepsilon^2}+i\frac\varepsilon{x^2+\varepsilon^2}$, et traitons les deux termes, ce qui est exactement (e).

\emph{Partie imaginaire.} $P_\varepsilon(x):=\frac1\pi\frac\varepsilon{x^2+\varepsilon^2}=\varepsilon^{-1}P(x/\varepsilon)$ avec $P(u)=\frac1{\pi(1+u^2)}$, positive d'intégrale $1$ ($\int\frac{du}{1+u^2}=\pi$). Pour $\varphi\in\D(\R)$, $\int P_\varepsilon\varphi=\int P(u)\varphi(\varepsilon u)\,du\to\varphi(0)\int P=\varphi(0)$ par convergence dominée ($|P(u)\varphi(\varepsilon u)|\le\norm\varphi_\infty P(u)\in L^1$). Donc $P_\varepsilon\wk\delta_0$, et $\frac\varepsilon{x^2+\varepsilon^2}\wk\pi\delta_0$. (C'est une approximation de l'unité au sens de la Définition~3.4, issue de $P\in L^1$ non à support compact ; l'Exercice~3.8 s'appliquerait, mais le calcul direct est plus simple.)

\emph{Partie réelle.} $Q_\varepsilon(x):=\frac x{x^2+\varepsilon^2}$ est impaire, donc $\int_{-R}^RQ_\varepsilon=0$, et pour $\supp\varphi\subset[-R,R]$ :
\[
\int Q_\varepsilon\varphi=\int_{-R}^RQ_\varepsilon(x)(\varphi(x)-\varphi(0))\,dx=\int_{-R}^R\frac{x^2}{x^2+\varepsilon^2}\cdot\frac{\varphi(x)-\varphi(0)}x\,dx .
\]
L'intégrande est majoré par $\big|\frac{\varphi(x)-\varphi(0)}x\big|\le\norm{\varphi'}_\infty$ (accroissements finis), intégrable sur $[-R,R]$, et converge simplement vers $\frac{\varphi(x)-\varphi(0)}x$ pour $x\ne0$. Par convergence dominée,
\[
\int Q_\varepsilon\varphi\longrightarrow\int_{-R}^R\frac{\varphi(x)-\varphi(0)}x\,dx=\langle\mathrm{vp}(1/x),\varphi\rangle
\]
(c'est la formule du polycopié avant la Définition~6.6 : $\mathrm{vp}\int\frac\varphi x=\int_{-R}^R\frac{\varphi(x)-\varphi(0)}x\,dx$, l'intégrande étant prolongeable en une fonction continue). Donc $Q_\varepsilon\wk\mathrm{vp}(1/x)$.

\emph{Conclusion.} $f_\varepsilon=Q_\varepsilon+i\pi P_\varepsilon\wk\mathrm{vp}(1/x)+i\pi\delta_0$. En conjuguant, $\frac1{x+i\varepsilon}=Q_\varepsilon-i\pi P_\varepsilon\wk\mathrm{vp}(1/x)-i\pi\delta_0$. Enfin $x\cdot\frac1{x-i\varepsilon}=1+\frac{i\varepsilon}{x-i\varepsilon}$, et $\big|\frac{i\varepsilon}{x-i\varepsilon}\big|\le1$ avec convergence simple vers $0$ pour $x\ne0$, donc $\frac{i\varepsilon}{x-i\varepsilon}\wk0$ (convergence dominée) ; la multiplication par $x$ étant continue sur $\D'$, $x\big(\mathrm{vp}(1/x)+i\pi\delta_0\big)=1$, et comme $x\delta_0=0$, $x\,\mathrm{vp}(1/x)=1$.

\medskip
\noindent\textbf{(e)} Déjà fait en (d).
\end{solution}

\begin{commentaire}
La formule de Sokhotski--Plemelj $\frac1{x\mp i0}=\mathrm{vp}\frac1x\pm i\pi\delta_0$ est le pont entre la valeur principale et l'analyse complexe (valeurs au bord de la fonction holomorphe $1/z$) ; elle est omniprésente en physique (fonctions de Green, relations de Kramers--Kronig). On y voit aussi que $\mathrm{vp}(1/x)$ n'est pas la seule « régularisation » naturelle de $1/x$ : $\frac1{x\mp i0}$ en sont deux autres, qui en diffèrent par un multiple de $\delta_0$ — exactement l'indétermination des solutions de $xT=1$. La partie finie $\mathrm{Pf}(1/x^2)$ est le second étage de la même construction (Hadamard) : soustraire à l'intégrale divergente son développement en $\varepsilon$, et ne garder que la partie finie ; ses solutions compagnes sont celles de $x^2T=0$.
\end{commentaire}

\section*{Convolution, solutions fondamentales, dimension supérieure}

% =====================================================================
\begin{exo}{Exercice 8 -- Convolution et solutions fondamentales}
Soit $\psi\in\D(\R)$.
\begin{enumerate}[label=(\alph*)]
  \item Calculer $\psi\star\delta_a$, $\psi\star\delta_0^{(k)}$, $\psi\star h$ et $\psi\star\mathbf 1_{[a,b]}$ (comme fonctions lisses, Proposition~6.10).
  \item Soit $E\in\D'(\R)$ une solution fondamentale de $P(\frac d{dx})=\sum_{j\le n}a_j\frac{d^j}{dx^j}$ ($a_j\in\C$), \ie\ $P(\frac d{dx})E=\delta_0$. Montrer que $u:=\psi\star E$ est une fonction lisse vérifiant $P(\frac d{dx})u=\psi$.
  \item Avec $E=he^{-x}$ (TD0 ch.~6) puis $E=h\sin x$ (Exercice~4), expliciter la solution $u$ de $u'+u=\psi$ puis de $u''+u=\psi$ ainsi obtenue. Quelle propriété supplémentaire a cette solution parmi toutes les solutions ?
  \item Montrer que si $T\in\D'(\R)$ vérifie $T'+T=0$ alors $T=ce^{-x}$, et en déduire que $u=\psi\star(he^{-x})$ est l'unique solution de $u'+u=\psi$ nulle au voisinage de $-\infty$.
  \item Soit $T\in\D'(\R)$ et $\psi\in\D(\R)$. Montrer que $(\psi\star T)'=\psi'\star T=\psi\star T'$, puis que $\supp(\psi\star T)\subset\supp\psi+\supp T$ (somme de Minkowski). Que peut-on dire de $\psi\star T$ lorsque $\supp T$ est compact ?
\end{enumerate}
\end{exo}

\begin{solution}
\noindent\textbf{(a)} Par la Proposition~6.10, $\psi\star T$ est la fonction lisse $x\mapsto\langle T,\psi(x-\cdot)\rangle$ : en effet, pour $\varphi\in\D$, $\langle\psi\star T,\varphi\rangle=\langle T,\varphi\star\check\psi\rangle$ avec $(\varphi\star\check\psi)(y)=\int\varphi(x)\psi(x-y)\,dx$, et l'on admet (formule classique, que la preuve de la Proposition~6.10 du polycopié n'établit pas : c'est le théorème de Fubini lorsque $T$ est une fonction, et en général on approche l'intégrale par des sommes de Riemann, qui convergent dans $\D(\R)$) que $\langle T,\int\varphi(x)\psi(x-\cdot)\,dx\rangle=\int\varphi(x)\langle T,\psi(x-\cdot)\rangle\,dx$. Pour les distributions ci-dessous on peut vérifier directement cette formule, et l'on obtient :
\begin{itemize}
  \item $\psi\star\delta_a=\psi(x-a)=\tau_a\psi$ : en effet $\langle\psi\star\delta_a,\varphi\rangle=(\varphi\star\check\psi)(a)=\int\varphi(x)\psi(x-a)\,dx$.
  \item $\psi\star\delta_0^{(k)}=\psi^{(k)}$ :
  \[
  \langle\delta_0^{(k)},\varphi\star\check\psi\rangle=(-1)^k(\varphi\star\check\psi)^{(k)}(0)=(-1)^k(\varphi\star\check\psi^{(k)})(0)=(-1)^k\int\varphi(x)\check\psi^{(k)}(-x)\,dx\ ;
  \]
  or $\check\psi^{(k)}(y)=(-1)^k\psi^{(k)}(-y)$, d'où $\int\varphi(x)\psi^{(k)}(x)\,dx$.
  \item $\psi\star h$ : $(\psi\star h)(x)=\int h(x-y)\psi(y)\,dy=\int_{-\infty}^x\psi(y)\,dy$, la primitive de $\psi$ nulle en $-\infty$. (Pour une fonction $T=T_f$, $\psi\star T_f=T_{\psi\star f}$, Exercice~6.4.)
  \item $\psi\star\mathbf 1_{[a,b]}=\int_{x-b}^{x-a}\psi(y)\,dy=\Psi(x-a)-\Psi(x-b)$ avec $\Psi$ la primitive ci-dessus : une « moyenne glissante ».
\end{itemize}

\medskip
\noindent\textbf{(b)} $u=\psi\star E$ est lisse (Proposition~6.10) et, par la même proposition, $\frac{d^j}{dx^j}(\psi\star E)=\psi\star E^{(j)}$ ; par linéarité, $P(\frac d{dx})u=\psi\star P(\frac d{dx})E=\psi\star\delta_0=\psi$ (Exemple~6.13).

\medskip
\noindent\textbf{(c)} $u(x)=(\psi\star he^{-x})(x)=\int h(x-y)e^{-(x-y)}\psi(y)\,dy=\int_{-\infty}^xe^{-(x-y)}\psi(y)\,dy$ : c'est la formule de variation de la constante avec « condition initiale en $-\infty$ ». De même $u(x)=\int_{-\infty}^x\sin(x-y)\psi(y)\,dy$ résout $u''+u=\psi$. Propriété supplémentaire : $u$ est nulle à gauche du support de $\psi$ ($u(x)=0$ si $x<\min\supp\psi$), ce que (e) exprime par $\supp u\subset\supp\psi+\R_+$. C'est la solution \emph{causale} : la réponse ne précède pas la source.

\medskip
\noindent\textbf{(d)} $(e^xT)'=e^x(T'+T)=0$ par Leibniz, donc $e^xT=c$ et $T=ce^{-x}$ (multiplier par $e^{-x}$). Si $u_1,u_2$ sont deux solutions de $u'+u=\psi$ nulles au voisinage de $-\infty$, $u_1-u_2=ce^{-x}$ est nulle au voisinage de $-\infty$, donc $c=0$.

\medskip
\noindent\textbf{(e)} $\langle(\psi\star T)',\varphi\rangle=-\langle T,\varphi'\star\check\psi\rangle$. Or $\varphi'\star\check\psi=(\varphi\star\check\psi)'=\varphi\star(\check\psi)'=-\varphi\star\check{(\psi')}$ (dérivée d'une convolution ; et $(\check\psi)'(y)=-\psi'(-y)$). Donc $\langle(\psi\star T)',\varphi\rangle=-\langle T,(\varphi\star\check\psi)'\rangle=\langle T',\varphi\star\check\psi\rangle=\langle\psi\star T',\varphi\rangle$, et aussi $=\langle T,\varphi\star\check{(\psi')}\rangle=\langle\psi'\star T,\varphi\rangle$.

\emph{Support.} L'ensemble $F:=\supp\psi+\supp T$ est fermé (somme d'un compact et d'un fermé) ; montrons que $\psi\star T$ est nulle sur l'ouvert $U:=\R\setminus F$. Soit $\varphi\in\D(U)$. On a $\supp(\varphi\star\check\psi)\subset\supp\varphi-\supp\psi$ (Exercice~6.4), et ce compact ne rencontre pas $\supp T$ : si $t=s-p$ avec $t\in\supp T$, $s\in\supp\varphi$ et $p\in\supp\psi$, alors $s=p+t\in F$, ce qui est exclu. Donc $\varphi\star\check\psi$ est une fonction test à support dans l'ouvert $\R\setminus\supp T$, où $T$ est nulle, et $\langle\psi\star T,\varphi\rangle=\langle T,\varphi\star\check\psi\rangle=0$. (Avec la formule ponctuelle admise en (a), cela se lit directement : pour $x\in U$, $y\mapsto\psi(x-y)$ a son support dans $x-\supp\psi$, disjoint de $\supp T$, donc $(\psi\star T)(x)=0$.) Si $\supp T$ est compact, $\supp\psi+\supp T$ est compact (image continue de $\supp\psi\times\supp T$), donc $\psi\star T$ est une fonction lisse à support compact : \emph{une fonction test}. C'est ce qui rend licite, dans la preuve du lemme de Weyl, l'écriture de $\psi\star(\eta S)$ comme fonction test. Attention : l'inclusion des supports peut être stricte, et $\psi\star T=0$ n'a aucune raison de se produire « parce que les supports sont loin » : $\psi\star\delta_0=\psi$ quel que soit $\supp\psi$.
\end{solution}

\begin{pointcle}
\textbf{$E\star\psi$ résout $P(\frac d{dx})u=\psi$.} C'est la raison d'être des solutions fondamentales : résoudre l'équation une fois pour toutes avec le second membre le plus concentré possible ($\delta_0$), puis superposer, la linéarité de l'équation et la commutation de $P(\frac d{dx})$ avec la convolution faisant le reste. Le même principe vaut pour le laplacien (Exercice~9) et pour toutes les équations à coefficients constants ; en dimension $1$ la solution fondamentale $h\cdot y$ est causale, ce qui sélectionne la solution « physique ».
\end{pointcle}

% =====================================================================
\begin{exo}{Exercice 9 -- Quelques calculs en dimension $d\ge2$}
\begin{enumerate}[label=(\alph*)]
  \item Dans $\R^2$, soit $Q:=\{x_1>0,\ x_2>0\}$ et $T:=\mathbf 1_Q$. Calculer $\partial_1T$ et $\partial_1\partial_2T$.
  \item Dans $\R^2$, soit $H:=\mathbf 1_{\{x_1>0\}}$. Calculer $\partial_1H$ et $\partial_2H$, et montrer que $\partial_1H$ n'est pas une fonction.
  \item Soit $E(x):=-\dfrac1{4\pi\norm x}$ sur $\R^3$. On sait que $\Delta E=\delta_0$ (Proposition~6.5). Pour $\psi\in\D(\R^3)$ on pose $u:=\psi\star E$, \ie\ $u(x)=-\dfrac1{4\pi}\displaystyle\int_{\R^3}\frac{\psi(y)}{\norm{x-y}}\,dy$ (potentiel newtonien). Montrer que $u\in\Cinf(\R^3)$, que $\Delta u=\psi$, et que $u(x)=O(1/\norm x)$ à l'infini.
  \item En déduire que si $v\in\D'(\R^3)$ vérifie $\Delta v=\psi$, alors $v-u$ est une fonction lisse harmonique ; et que $u$ est l'unique solution de $\Delta v=\psi$ qui soit une fonction lisse tendant vers $0$ à l'infini et dont les dérivées premières sont dans $L^2(\R^3)$. \emph{On utilisera le Corollaire~6.3.}
\end{enumerate}
\end{exo}

\begin{solution}
\noindent\textbf{(a)} Pour $\varphi\in\D(\R^2)$ : $\langle\partial_1T,\varphi\rangle=-\int_Q\partial_1\varphi=-\int_0^\infty\Big(\int_0^\infty\partial_1\varphi(x_1,x_2)\,dx_1\Big)dx_2=\int_0^\infty\varphi(0,x_2)\,dx_2$. Donc $\partial_1T$ est la mesure « longueur sur la demi-droite $\{0\}\times\R_+$ » (comparer avec l'Exemple~6.7). Puis $\langle\partial_1\partial_2T,\varphi\rangle=-\langle\partial_1T,\partial_2\varphi\rangle=-\int_0^\infty\partial_2\varphi(0,x_2)\,dx_2=\varphi(0,0)$ : $\partial_1\partial_2\mathbf 1_Q=\delta_0$. C'est la version bidimensionnelle de $h'=\delta_0$ : $\mathbf 1_Q=h(x_1)h(x_2)$ et $\partial_1\partial_2(h\otimes h)=h'\otimes h'=\delta_0\otimes\delta_0$.

\medskip
\noindent\textbf{(b)} $\langle\partial_1H,\varphi\rangle=-\int_{x_1>0}\partial_1\varphi=\int_\R\varphi(0,x_2)\,dx_2=:\langle\Lambda,\varphi\rangle$, la mesure de longueur sur l'axe $\{x_1=0\}$ ; et $\langle\partial_2H,\varphi\rangle=-\int_{x_1>0}\partial_2\varphi=-\int_0^\infty\big[\varphi(x_1,\cdot)\big]_{-\infty}^{+\infty}dx_1=0$. \emph{$\Lambda$ n'est pas une fonction :} même argument que pour l'Exemple~6.7 (TD0 ch.~6) : $\Lambda$ est nulle sur l'ouvert $\{x_1\ne0\}$ ; si $\Lambda=T_g$, le lemme de du Bois-Reymond sur cet ouvert donne $g=0$ p.p. sur $\{x_1\ne0\}$, donc p.p. sur $\R^2$ (l'axe est négligeable), donc $\Lambda=0$ ; mais $\Lambda(\varphi)>0$ pour $\varphi\ge0$ non nulle sur l'axe.

\medskip
\noindent\textbf{(c)} $E\in L^1_{\mathrm{loc}}(\R^3)$ (Proposition~6.5), donc $\psi\star E=\psi\star T_E=T_{\psi\star E}$ avec la convolution usuelle $L^1_{\mathrm{loc}}\star\D$ (Exercice~6.4), et $u(x)=\int\psi(x-y)E(y)\,dy$. Régularité : par dérivation sous l'intégrale (TDS : $y\mapsto\partial^\alpha\psi(x-y)E(y)$ est dominée, pour $x$ dans un borné $B$, par $\norm{\partial^\alpha\psi}_\infty|E|\mathbf 1_{B-\supp\psi}\in L^1$), $u\in\Cinf$ avec $\partial^\alpha u=(\partial^\alpha\psi)\star E$ ; c'est aussi la Proposition~6.10. Alors $\Delta u=(\Delta\psi)\star E=\psi\star\Delta E=\psi\star\delta_0=\psi$ (Proposition~6.10 et Exemple~6.13). Comportement à l'infini : si $\supp\psi\subset\Bf(0,R)$ et $\norm x\ge2R$, alors pour $y\in\supp\psi$, $\norm{x-y}\ge\norm x-R\ge\norm x/2$, donc $|u(x)|\le\frac1{4\pi}\norm\psi_1\cdot\frac2{\norm x}=O(1/\norm x)$. (De même $\nabla u=O(\norm x^{-2})$ en dérivant $1/\norm{x-y}$, d'où $\nabla u\in L^2(\R^3)$ : $\int_{\norm x\ge2R}\norm x^{-4}\,dx<\infty$ en dimension $3$.)

\medskip
\noindent\textbf{(d)} $\Delta(v-u)=\psi-\psi=0$ : $v-u$ est une distribution harmonique, donc une fonction lisse par le lemme de Weyl (Lemme~6.4). Si de plus $v$ est lisse, tend vers $0$ à l'infini et $\nabla v\in L^2$, alors $w:=v-u$ est lisse harmonique, tend vers $0$ à l'infini (car $u$ aussi) et $\nabla w\in L^2$ (car $\nabla u\in L^2$ par (c)). Le Corollaire~6.3 (Liouville) donne $w$ constante, et la limite nulle à l'infini force $w=0$ : $v=u$.
\end{solution}

\begin{commentaire}
Le point (c) est le \emph{potentiel newtonien} : $u$ est le potentiel gravitationnel (ou électrostatique) créé par la densité $\psi$, et $\Delta u=\psi$ est l'équation de Poisson. Le point (d) montre que la décroissance à l'infini joue le rôle de « condition au bord » qui sélectionne la solution physique, exactement comme la causalité en dimension $1$ (Exercice~8). Le lemme de Weyl et le théorème de Liouville sont les deux énoncés de rigidité qui rendent la théorie efficace : chercher une solution dans l'énorme espace $\D'$ ne fait pas apparaître de solutions parasites, seulement des solutions lisses.
\end{commentaire}

\section*{Distributions périodiques et séries de Fourier}

% =====================================================================
\begin{exo}{Exercice 10 -- Distributions périodiques}
On dit que $T\in\D'(\R)$ est \emph{$2\pi$-périodique} si $\tau_{2\pi}T=T$. On fixe $\theta\in\D(\R)$ telle que $\sum_{n\in\Z}\theta(x+2\pi n)=1$ pour tout $x\in\R$ (une telle fonction existe : prendre $\theta:=\theta_0/\Theta_0$ avec $\theta_0$ une bosse positive valant $1$ sur $[0,2\pi]$ et $\Theta_0:=\sum_n\theta_0(\cdot+2\pi n)$, lisse, $2\pi$-périodique et $\ge1$).
\begin{enumerate}[label=(\alph*)]
  \item Montrer que $T_f$ est $2\pi$-périodique si $f\in L^1_{\mathrm{loc}}$ est $2\pi$-périodique, et que le peigne $\Sha:=\sum_{n\in\Z}\delta_{2\pi n}$ est $2\pi$-périodique. Montrer que les dérivées d'une distribution périodique sont périodiques.
  \item Soit $\psi\in\D(\R)$ avec $\sum_n\psi(x+2\pi n)=0$ pour tout $x$. Montrer que $\psi=\tau_{2\pi}\chi-\chi$ pour une $\chi\in\D(\R)$. \emph{Indication : $\chi:=\sum_{n\ge1}\psi(\cdot+2\pi n)$.} En déduire que si $T$ est $2\pi$-périodique, $\langle T,\psi g\rangle=0$ pour toute fonction $g\in\Cinf(\R)$ $2\pi$-périodique.
  \item Pour $T$ $2\pi$-périodique et $k\in\Z$, on pose $c_k(T):=\frac1{2\pi}\langle T,\theta\,e^{-ikx}\rangle$. Montrer que $c_k(T)$ ne dépend pas du choix de $\theta$, et que pour $T=T_f$ on retrouve $c_k(f)=\frac1{2\pi}\int_0^{2\pi}f(x)e^{-ikx}\,dx$. Calculer $c_k(\Sha)$.
  \item Montrer que $c_k(T')=ik\,c_k(T)$.
  \item Montrer qu'il existe $m\in\N$ et $C>0$ tels que $|c_k(T)|\le C(1+|k|)^m$ pour tout $k\in\Z$ (croissance polynomiale).
  \item En déduire que la série $\sum_{k\in\Z}c_k(T)e^{ikx}$ converge dans $\D'(\R)$ (TD0 ch.~6, séries trigonométriques). On admet que sa somme est $T$ (théorème de Fourier pour les distributions périodiques). Vérifier ce résultat sur $\Sha$ à l'aide du TD0 ch.~6 (peigne de Dirac), et sur $T=\Sha'$.
\end{enumerate}
\end{exo}

\begin{solution}
\noindent\textbf{(a)} $\tau_{2\pi}T_f=T_{\tau_{2\pi}f}$ (Exercice~6.3) et $\tau_{2\pi}f=f(\cdot-2\pi)=f$. Pour le peigne : $\langle\tau_{2\pi}\Sha,\varphi\rangle=\langle\Sha,\varphi(\cdot+2\pi)\rangle=\sum_n\varphi(2\pi n+2\pi)=\sum_n\varphi(2\pi(n+1))=\langle\Sha,\varphi\rangle$ (réindexation d'une somme finie). Dérivées : $\tau_a$ et $\partial$ commutent (Exercice~6.3), donc $\tau_{2\pi}T'=(\tau_{2\pi}T)'=T'$.

\medskip
\noindent\textbf{(b)} Soit $\supp\psi\subset[-R,R]$ et $\chi(x):=\sum_{n\ge1}\psi(x+2\pi n)$. Pour $x$ fixé, seuls les $n$ avec $|x+2\pi n|\le R$ contribuent : somme localement finie, $\chi$ est lisse. Si $x>R-2\pi$, alors $x+2\pi n>R$ pour tout $n\ge1$ : $\chi(x)=0$. Si $x<-R$, l'hypothèse $\sum_{n\in\Z}\psi(x+2\pi n)=0$ donne $\chi(x)=-\sum_{n\le0}\psi(x+2\pi n)$, et $x+2\pi n\le x<-R$ pour $n\le0$ : $\chi(x)=0$. Donc $\supp\chi\subset[-R,R-2\pi]$, compact : $\chi\in\D(\R)$. (Si $R<\pi$ ce segment est vide, donc $\chi=0$ et $\psi=0$ : en effet $\psi$ et ses translatées ont alors des supports disjoints et l'hypothèse force $\psi=0$.) Enfin $\tau_{2\pi}\chi(x)-\chi(x)=\chi(x-2\pi)-\chi(x)=\sum_{n\ge1}\psi(x+2\pi(n-1))-\sum_{n\ge1}\psi(x+2\pi n)=\psi(x)$ (télescopage : $\sum_{n\ge0}-\sum_{n\ge1}$).

Soit alors $T$ périodique et $g$ lisse $2\pi$-périodique. $\psi g=(\tau_{2\pi}\chi-\chi)g=\tau_{2\pi}(\chi g)-\chi g$ car $g=\tau_{2\pi}g$. Donc $\langle T,\psi g\rangle=\langle T,\tau_{2\pi}(\chi g)\rangle-\langle T,\chi g\rangle=\langle\tau_{-2\pi}T,\chi g\rangle-\langle T,\chi g\rangle=0$, puisque $\tau_{2\pi}T=T$ équivaut à $\tau_{-2\pi}T=T$ (appliquer $\tau_{-2\pi}$, et $\tau_a\tau_b=\tau_{a+b}$).

\medskip
\noindent\textbf{(c)} Si $\theta_1$ est un autre choix, $\psi:=\theta-\theta_1$ vérifie $\sum_n\psi(x+2\pi n)=1-1=0$, et $e^{-ikx}$ est lisse $2\pi$-périodique : par (b), $\langle T,(\theta-\theta_1)e^{-ikx}\rangle=0$. Pour $T=T_f$ : $\langle T_f,\theta e^{-ikx}\rangle=\int_\R f\theta e^{-ikx}=\sum_n\int_{2\pi n}^{2\pi(n+1)}f(x)\theta(x)e^{-ikx}\,dx=\sum_n\int_0^{2\pi}f(y)\theta(y+2\pi n)e^{-iky}\,dy$ (changement $x=y+2\pi n$, périodicité de $f$ et de $e^{-ikx}$) $=\int_0^{2\pi}f(y)e^{-iky}\Big(\sum_n\theta(y+2\pi n)\Big)dy=\int_0^{2\pi}fe^{-iky}$ (Fubini licite, sommes finies sur $[0,2\pi]$). Pour le peigne : $\langle\Sha,\theta e^{-ikx}\rangle=\sum_n\theta(2\pi n)e^{-2i\pi kn}=\sum_n\theta(2\pi n)=1$, donc $c_k(\Sha)=\frac1{2\pi}$ pour tout $k$.

\medskip
\noindent\textbf{(d)} $c_k(T')=\frac1{2\pi}\langle T',\theta e^{-ikx}\rangle=-\frac1{2\pi}\langle T,(\theta e^{-ikx})'\rangle=-\frac1{2\pi}\langle T,\theta'e^{-ikx}\rangle+\frac{ik}{2\pi}\langle T,\theta e^{-ikx}\rangle$. Le premier terme est nul par (b) : $\sum_n\theta'(x+2\pi n)=\frac d{dx}\sum_n\theta(x+2\pi n)=0$. Donc $c_k(T')=ik\,c_k(T)$.

\medskip
\noindent\textbf{(e)} Soit $K:=\supp\theta$, et $C_K,m_K$ les constantes de (6.1) pour $T$. Par Leibniz, $(\theta e^{-ikx})^{(j)}=\sum_{i\le j}\binom ji\theta^{(j-i)}(-ik)^ie^{-ikx}$, donc $\norm{\theta e^{-ikx}}_{m_K,K}\le C'(1+|k|)^{m_K}$. D'où $|c_k(T)|\le\frac{C_KC'}{2\pi}(1+|k|)^{m_K}$.

\medskip
\noindent\textbf{(f)} Par (e), $|c_k(T)|\lesssim(1+|k|)^m$ ; la série $\sum_kc_k(T)(ik)^{-m-2}e^{ikx}$ (termes $k\ne0$) converge normalement, vers une fonction continue $F$, et
\[
\sum_{|k|\le N}c_k(T)e^{ikx}=c_0(T)+\Big(\sum_{0<|k|\le N}c_k(T)(ik)^{-m-2}e^{ikx}\Big)^{(m+2)}\wk c_0(T)+F^{(m+2)}
\]
par continuité de la dérivation (TD0 ch.~6, Exemple~6.2). Pour $\Sha$ : $c_k(\Sha)=\frac1{2\pi}$ et le TD0 (peigne de Dirac) donne exactement $\sum_ke^{ikx}=2\pi\Sha$, soit $\Sha=\sum_k\frac1{2\pi}e^{ikx}$ : cohérent. Pour $\Sha'$ : $c_k(\Sha')=\frac{ik}{2\pi}$ par (d), et $\sum_k\frac{ik}{2\pi}e^{ikx}=\big(\sum_k\frac1{2\pi}e^{ikx}\big)'=\Sha'$ par continuité de la dérivation : cohérent aussi. (La preuve générale de « la somme est $T$ » consiste à montrer que $T-\sum c_ke^{ikx}$ est périodique à coefficients nuls, puis qu'une distribution périodique à coefficients de Fourier tous nuls est nulle, par densité des polynômes trigonométriques dans les fonctions lisses périodiques pour toutes les normes $\mathcal C^m$ — par exemple via les sommes de Fejér, qui sont des approximations de l'unité périodiques.)
\end{solution}

\begin{pointcle}
\textbf{Les distributions périodiques sont exactement les séries trigonométriques à coefficients à croissance polynomiale.} Le sens « coefficients polynomiaux $\Rightarrow$ convergence dans $\D'$ » est la régularisation par intégration du TD0 ; le sens « distribution $\Rightarrow$ coefficients polynomiaux » est l'estimation (6.1), l'ordre de $T$ donnant la puissance. Entre les deux, la formule $c_k(T')=ik\,c_k(T)$ : dériver, c'est multiplier par $ik$, et une distribution d'ordre $m$ a des coefficients en $O(k^m)$. Les fonctions $L^2$ sont le cas $c_k\in\ell^2$ (chapitre~4), les fonctions lisses le cas « décroissance rapide », et le peigne de Dirac, avec ses coefficients constants, est la distribution périodique « la plus singulière à l'ordre $0$ ». C'est le point de départ de la transformée de Fourier des distributions tempérées, hors programme mais directement dans la continuité de ce cours.
\end{pointcle}

\section*{Vers l'examen : questions de cours, problème transversal, compléments}

\begin{commentaire}
Ces quatre exercices complètent la feuille dans l'optique de l'examen. L'Exercice~11 fait restituer les preuves du polycopié ; l'Exercice~12 est un problème long qui relie les chapitres 4, 5 et 6 autour de la Proposition~6.6 ; l'Exercice~13 traite deux points du polycopié absents des exercices précédents (la convexité de l'Exemple~6.11 et un opérateur non hypoelliptique, à comparer à la Remarque~6.9) ; l'Exercice~14 est un vrai ou faux sur les pièges du chapitre.
\end{commentaire}

% =====================================================================
\begin{exo}{Exercice 11 -- Questions de cours}
\emph{À traiter sans document, en rédigeant comme le jour de l'examen : énoncé complet, puis démonstration.}
\begin{enumerate}[label=(\alph*)]
  \item Énoncer et démontrer le Théorème~6.2 (une forme linéaire positive sur $\D(\Omega)$ est une distribution d'ordre $0$).
  \item Énoncer et démontrer le lemme de du Bois-Reymond (Lemme~6.2).
  \item Énoncer la formule des sauts (Théorème~6.3) et la démontrer lorsque $A=\{a\}$ est réduit à un point.
  \item Soit $I\subset\R$ un intervalle ouvert et $T\in\D'(I)$ telle que $T'=0$. Démontrer que $T$ est constante (Théorème~6.4 en dimension~$1$ ; on redémontrera le cas utile du Lemme~6.3).
  \item Énoncer et démontrer la Proposition~6.2 ($|T(\varphi)|\lesssim\norm\varphi_p$ entraîne $T\in L^{p'}(\Omega)$, pour $p<\infty$). Que se passe-t-il pour $p=\infty$ ?
\end{enumerate}
\end{exo}

\begin{solution}
\noindent\textbf{(a)} \emph{Énoncé.} Soit $T$ une forme linéaire sur $\D(\Omega)$, positive au sens où $\varphi\ge0\Rightarrow T(\varphi)\ge0$. Alors $T$ est une distribution d'ordre $0$.

\emph{Preuve.} Soit $K\subset\Omega$ compact et $\theta_K\in\D(\Omega)$ à valeurs dans $[0,1]$, valant $1$ sur $K$ (Lemme~6.1). Soit d'abord $\varphi\in\D_K(\Omega)$ \emph{réelle}. Comme $|\varphi|\le\norm\varphi_K$ et que $\varphi$ est nulle hors de $K$, où $\theta_K=1$, on a partout
\[
-\norm\varphi_K\,\theta_K\le\varphi\le\norm\varphi_K\,\theta_K .
\]
Les fonctions $\norm\varphi_K\theta_K\pm\varphi$ sont donc des fonctions tests positives, et la positivité et la linéarité de $T$ donnent $\pm T(\varphi)\le\norm\varphi_KT(\theta_K)$, soit $|T(\varphi)|\le T(\theta_K)\norm\varphi_K$ (en particulier $T(\varphi)$ est réel, et $T(\theta_K)\ge0$). Si $\varphi=u+iv\in\D_K(\Omega)$ est complexe, $u$ et $v$ sont dans $\D_K(\Omega)$, réelles, de norme $\le\norm\varphi_K$, et $|T(\varphi)|\le|T(u)|+|T(v)|\le2T(\theta_K)\norm\varphi_K$. C'est l'estimation (6.1) avec $m_K=0$ et $C_K=2T(\theta_K)$.

\medskip
\noindent\textbf{(b)} \emph{Énoncé.} L'application $f\mapsto T_f$ de $L^1_{\mathrm{loc}}(\Omega)$ dans $\D'(\Omega)$ est injective : si $\int_\Omega f\varphi=0$ pour toute $\varphi\in\D(\Omega)$, alors $f=0$ presque partout.

\emph{Preuve.} Soit $f\in L^1_{\mathrm{loc}}(\Omega)$ avec $T_f=0$, et $\theta\in\D(\Omega)$. La fonction $g:=\theta f$, prolongée par $0$ hors de $\Omega$, est dans $L^1(\R^d)$ ($f$ est intégrable sur le compact $\supp\theta$). Soit $(\rho_n)$ une approximation de l'unité issue d'une fonction test à support dans la boule unité. Pour tout $x\in\R^d$, la fonction $y\mapsto\theta(y)\rho_n(x-y)$ est une fonction test sur $\Omega$ (lisse, à support dans $\supp\theta$), donc
\[
(g\star\rho_n)(x)=\int_\Omega f(y)\,\theta(y)\rho_n(x-y)\,dy=\big\langle T_f,\theta\,\rho_n(x-\cdot)\big\rangle=0 .
\]
Or $g\star\rho_n\to g$ dans $L^1(\R^d)$ (Théorème~3.6, $p=1$) : donc $g=0$ dans $L^1$, \ie\ $\theta f=0$ presque partout. En prenant $\theta=\theta_n$ valant $1$ sur le compact $K_n$ d'une suite exhaustive (Lemme~6.1), $f=0$ p.p.\ sur chaque $K_n$, donc sur $\Omega=\bigcup_nK_n$ (réunion dénombrable d'ensembles négligeables).

\medskip
\noindent\textbf{(c)} \emph{Énoncé.} Soit $A\subset\R$ fini et $f$ de classe $\mathcal C^1$ sur $\R\setminus A$, telle que $f$ et $f'$ admettent des limites finies à droite et à gauche en tout point de $A$. Alors $f$ et $f'$ sont dans $L^1_{\mathrm{loc}}(\R)$ et $T_f'=T_{f'}+\sum_{a\in A}\big(f(a^+)-f(a^-)\big)\delta_a$.

\emph{Preuve pour $A=\{a\}$.} $f$ et $f'$ sont continues sur $\R\setminus\{a\}$ et bornées au voisinage de $a$ (limites finies) : elles sont localement bornées, donc dans $L^1_{\mathrm{loc}}$. La fonction $f_g$ égale à $f$ sur $]-\infty,a[$ et à $f(a^-)$ en $a$ est continue sur $]-\infty,a]$, dérivable sur $]-\infty,a[$ avec $f_g'=f'\to f'(a^-)$ : par le théorème de la limite de la dérivée, $f_g$ est $\mathcal C^1$ sur $]-\infty,a]$. De même $f_d$ ($f$ prolongée par $f(a^+)$) est $\mathcal C^1$ sur $[a,+\infty[$. Soit $\varphi\in\D(\R)$, à support dans $[-R,R]$ avec $R>|a|$. On intègre par parties sur chaque morceau :
\begin{align*}
\int_{-R}^af\varphi'&=\big[f_g\varphi\big]_{-R}^a-\int_{-R}^af'\varphi=f(a^-)\varphi(a)-\int_{-R}^af'\varphi,\\
\int_a^Rf\varphi'&=\big[f_d\varphi\big]_a^R-\int_a^Rf'\varphi=-f(a^+)\varphi(a)-\int_a^Rf'\varphi .
\end{align*}
En sommant, $\langle T_f',\varphi\rangle=-\int_\R f\varphi'=\int_\R f'\varphi+\big(f(a^+)-f(a^-)\big)\varphi(a)$, ce qui est la formule annoncée.

\medskip
\noindent\textbf{(d)} \emph{Étape 1 : une fonction test d'intégrale nulle est la dérivée d'une fonction test.} Soit $\psi\in\D(I)$ avec $\int_I\psi=0$, et $[\alpha,\beta]\subset I$ un segment contenant $\supp\psi$. Posons $\Psi(x):=\int_\alpha^x\psi$ pour $x\in I$. Alors $\Psi$ est lisse, $\Psi'=\psi$, $\Psi=0$ sur $I\cap]-\infty,\alpha]$, et pour $x\ge\beta$, $\Psi(x)=\int_I\psi=0$. Donc $\supp\Psi\subset[\alpha,\beta]$ : $\Psi\in\D(I)$.

\emph{Étape 2.} Fixons $\theta\in\D(I)$ d'intégrale $1$ et posons $c:=T(\theta)$. Pour $\varphi\in\D(I)$, la fonction $\psi:=\varphi-\big(\int_I\varphi\big)\theta$ est une fonction test d'intégrale nulle : par l'étape~1, $\psi=\Psi'$ avec $\Psi\in\D(I)$, d'où $T(\psi)=T(\Psi')=-T'(\Psi)=0$. Par linéarité,
\[
T(\varphi)=\Big(\int_I\varphi\Big)T(\theta)=c\int_I\varphi=\langle T_c,\varphi\rangle ,
\]
\ie\ $T$ est la (distribution associée à la) fonction constante $c$. On remarquera que la connexité de $I$ sert à l'étape~1 : sur $I=]0,1[\,\cup\,]2,3[$, une fonction test d'intégrale nulle n'est pas nécessairement une dérivée de fonction test, et $\mathbf 1_{]0,1[}$ est de dérivée nulle sans être constante.

\medskip
\noindent\textbf{(e)} \emph{Énoncé.} Soit $p\in[1,\infty[$ et $T\in\D'(\Omega)$. S'il existe $C$ telle que $|T(\varphi)|\le C\norm\varphi_p$ pour toute $\varphi\in\D(\Omega)$, alors il existe $f\in L^{p'}(\Omega)$ telle que $T=T_f$.

\emph{Preuve.} $\D(\Omega)$ est un sous-espace dense de $L^p(\Omega)$ puisque $p<\infty$ (Théorème~6.1). L'hypothèse dit que $T$, linéaire, est continue sur $\D(\Omega)$ pour la norme $\norm\cdot_p$, donc lipschitzienne et uniformément continue. Comme $\C$ est complet, le Théorème~1.1 fournit un unique prolongement continu $\tilde T:L^p(\Omega)\to\C$ ; il est linéaire (passage à la limite dans $T(\lambda\varphi_n+\psi_n)=\lambda T(\varphi_n)+T(\psi_n)$) et vérifie $|\tilde T(g)|\le C\norm g_p$ : $\tilde T\in L^p(\Omega)'$. Par le Théorème~5.6 de représentation de Riesz ($1\le p<\infty$), il existe $f\in L^{p'}(\Omega)$ avec $\tilde T(g)=\int_\Omega fg$ pour tout $g\in L^p(\Omega)$. En particulier pour $g=\varphi\in\D(\Omega)$ : $T(\varphi)=\int f\varphi$, \ie\ $T=T_f$ (et $f\in L^{p'}\subset L^1_{\mathrm{loc}}$).

\emph{Cas $p=\infty$.} L'énoncé est faux : $\delta_a$ vérifie $|\delta_a(\varphi)|\le\norm\varphi_\infty$ et n'est pas une fonction (Exemple~6.4). Les deux ingrédients de la preuve tombent : $\D(\Omega)$ n'est pas dense dans $L^\infty(\Omega)$, et le dual de $L^\infty$ n'est pas $L^1$. Ce que l'on obtient est une distribution d'ordre $0$, c'est-à-dire une mesure de Radon (Remarque~6.5).
\end{solution}

\begin{methode}
\textbf{Rédiger une question de cours.} (1) \emph{L'énoncé d'abord}, avec toutes ses hypothèses ($p<\infty$, $A$ fini, $\Omega$ connexe…) : une hypothèse oubliée coûte plus cher qu'une maladresse dans la preuve. (2) Nommer chaque résultat utilisé au moment où il sert : Lemme~6.1 (plateau), Théorème~3.6 (approximation de l'unité), Théorème~1.1 (prolongement), Théorème~5.6 (Riesz). (3) Repérer le \emph{geste} de chaque preuve, qui tient en une ligne : encadrer $\varphi$ par $\pm\norm\varphi_K\theta_K$ pour (a) ; tester contre $\theta\rho_n(x-\cdot)$ pour (b) ; intégrer par parties morceau par morceau pour (c) ; corriger $\varphi$ par $(\int\varphi)\theta$ pour (d) ; prolonger par densité puis représenter pour (e). Ces cinq gestes, avec la dilatation d'une fonction test (ordre de $\delta_0^{(k)}$), couvrent presque toutes les preuves du chapitre.
\end{methode}


% =====================================================================
\begin{exo}{Exercice 12 -- Problème : séries de Fourier, dérivée dans $L^2$ et résolution spectrale (chapitres 4, 5 et 6)}
On note $J:=\,]0,2\pi[$ et $H:=L^2(J)$ muni du produit scalaire $\langle f,g\rangle:=\frac1{2\pi}\int_Jf\bar g$, pour lequel $e_k:x\mapsto e^{ikx}$ ($k\in\Z$) est une base hilbertienne (Exemple~4.7) ; $c_k(f):=\langle f,e_k\rangle$ et $S_Nf:=\sum_{|k|\le N}c_k(f)e_k$. Pour $f\in H$ on note $f^\sharp$ son prolongement $2\pi$-périodique à $\R$ (défini presque partout), qui est dans $L^2_{\mathrm{loc}}(\R)$, et l'on considère $S_Nf$ comme une fonction sur $\R$. On pose enfin
\[
H^1_\sharp:=\Big\{u\in H:\ \sum_{k\in\Z}k^2|c_k(u)|^2<\infty\Big\},\qquad L^2_0:=\{f\in H:\ c_0(f)=0\}.
\]
\begin{enumerate}[label=(\alph*)]
  \item Soit $f\in H$. Montrer que $S_Nf\to f^\sharp$ dans $L^2([a,b])$ pour tout segment $[a,b]$, puis dans $\D'(\R)$. En déduire que $\sum_{|k|\le N}ik\,c_k(f)e_k\wk(f^\sharp)'$ dans $\D'(\R)$.
  \item Soit $u\in H^1_\sharp$. Montrer qu'il existe $v\in H$ tel que $(u^\sharp)'=v^\sharp$ dans $\D'(\R)$ ; on note $u':=v$. Montrer que la série de Fourier de $u$ converge normalement sur $\R$, que $u$ admet un représentant continu $\tilde u$ sur $[0,2\pi]$ avec $\tilde u(0)=\tilde u(2\pi)$, et que
  \[
  \norm{\tilde u-c_0(u)}_\infty\le\frac\pi{\sqrt3}\,\norm{u'}_H .
  \]
  Réciproquement, montrer à l'aide de l'Exercice~10 que si $u,w\in H$ vérifient $(u^\sharp)'=w^\sharp$ dans $\D'(\R)$, alors $u\in H^1_\sharp$.
  \item Soit $f(x):=x$ sur $J$. Calculer $c_k(f)$. Vérifier que $T_f'=T_1$ dans $\D'(J)$, mais que $f\notin H^1_\sharp$. Calculer $(f^\sharp)'$ dans $\D'(\R)$ de deux façons et commenter.
  \item (Proposition~6.6.) Soit $f\in L^2_0$. Montrer qu'il existe un unique $\varphi_f\in L^2_0$ tel que $(\varphi_f^\sharp)''=f^\sharp$ dans $\D'(\R)$, et l'expliciter. Montrer que $\varphi_f$ admet un représentant de classe $\mathcal C^1$. Montrer enfin que l'unicité est \emph{fausse} si l'on demande seulement $\varphi''=f$ dans $\D'(J)$.
  \item Montrer que $K:f\mapsto\varphi_f$ est un opérateur linéaire continu, auto-adjoint et compact de l'espace de Hilbert $L^2_0$ ; calculer $\opnorm K$, ses valeurs propres et ses espaces propres, et comparer avec la Proposition~4.7 et le Théorème~5.5.
  \item On munit $H^1_\sharp$ de $\langle u,v\rangle_1:=\langle u,v\rangle+\langle u',v'\rangle$. Montrer que c'est un espace de Hilbert séparable. Soit $f\in H$ ; à l'aide du Théorème~4.6 de Riesz, montrer qu'il existe un unique $u\in H^1_\sharp$ tel que $\langle v,u\rangle_1=\langle v,f\rangle$ pour tout $v\in H^1_\sharp$, calculer ses coefficients de Fourier et vérifier que $-(u^\sharp)''+u^\sharp=f^\sharp$ dans $\D'(\R)$.
\end{enumerate}
\end{exo}

\begin{solution}
On utilisera la forme polarisée de l'identité de Parseval (Théorème~4.5 et continuité du produit scalaire),
\[
\langle f,g\rangle=\sum_{k\in\Z}c_k(f)\overline{c_k(g)},
\]
et le fait que $c_k\big(\sum_ja_je_j\big)=a_k$ pour toute série convergente dans $H$.

\medskip
\noindent\textbf{(a)} Par le Théorème~4.5, $\norm{S_Nf-f}_H\to0$. La fonction $S_Nf-f^\sharp$ est $2\pi$-périodique ; un segment $[a,b]$ est contenu dans $[-2\pi m,2\pi m]$ pour un $m\in\N^\star$, réunion de $2m$ périodes, donc
\[
\int_a^b|S_Nf-f^\sharp|^2\le2m\int_0^{2\pi}|S_Nf-f|^2=4\pi m\,\norm{S_Nf-f}_H^2\xrightarrow[N\to\infty]{}0 .
\]
Par Cauchy--Schwarz la convergence a lieu aussi dans $L^1([a,b])$ : $S_Nf\to f^\sharp$ dans $L^1_{\mathrm{loc}}(\R)$, donc dans $\D'(\R)$ (Exemple~6.1). La dérivation étant continue sur $\D'(\R)$ (Proposition~6.3) et $S_Nf$ étant lisse de dérivée classique $\sum_{|k|\le N}ik\,c_k(f)e_k$, on obtient $\sum_{|k|\le N}ik\,c_k(f)e_k=(S_Nf)'\wk(f^\sharp)'$.

\medskip
\noindent\textbf{(b)} \emph{Existence de $u'$.} Les vecteurs $ik\,c_k(u)e_k$ sont deux à deux orthogonaux et $\sum_kk^2|c_k(u)|^2<\infty$ : par la Proposition~4.4, la série $v:=\sum_kik\,c_k(u)e_k$ converge dans $H$, et $c_k(v)=ik\,c_k(u)$. Ses sommes partielles sont $S_Nv=(S_Nu)'$. Par (a) appliqué à $v$, $S_Nv\wk v^\sharp$ dans $\D'(\R)$ ; par (a) appliqué à $u$, $(S_Nu)'\wk(u^\sharp)'$. Par unicité de la limite dans $\D'(\R)$, $(u^\sharp)'=v^\sharp$.

\emph{Convergence normale.} Par l'inégalité de Cauchy--Schwarz dans $\ell^2(\Z^\star)$ et $\sum_{k\ne0}k^{-2}=2\cdot\frac{\pi^2}6=\frac{\pi^2}3$,
\[
\sum_{k\ne0}|c_k(u)|=\sum_{k\ne0}\frac1{|k|}\cdot|k|\,|c_k(u)|\le\Big(\sum_{k\ne0}\frac1{k^2}\Big)^{1/2}\Big(\sum_{k}k^2|c_k(u)|^2\Big)^{1/2}=\frac\pi{\sqrt3}\,\norm{u'}_H ,
\]
la dernière égalité par Parseval pour $u'=v$. Comme $\norm{e_k}_\infty=1$, la série $\sum_kc_k(u)e_k$ converge normalement sur $\R$ vers une fonction $\tilde u$ continue et $2\pi$-périodique, en particulier $\tilde u(0)=\tilde u(2\pi)$. Elle converge aussi vers $u$ dans $H$ ; la convergence uniforme sur $[0,2\pi]$ entraînant la convergence dans $H$, $\tilde u=u$ presque partout sur $J$. Enfin $\tilde u-c_0(u)=\sum_{k\ne0}c_k(u)e_k$, d'où $\norm{\tilde u-c_0(u)}_\infty\le\sum_{k\ne0}|c_k(u)|\le\frac\pi{\sqrt3}\norm{u'}_H$.

\emph{Réciproque.} Si $(u^\sharp)'=w^\sharp$, les distributions $T_{u^\sharp}$ et $T_{w^\sharp}$ sont $2\pi$-périodiques, et l'Exercice~10 donne $c_k(T_{w^\sharp})=c_k\big((T_{u^\sharp})'\big)=ik\,c_k(T_{u^\sharp})$ (question (d)), avec $c_k(T_{g^\sharp})=c_k(g)$ pour $g\in H$ (question (c)). Donc $c_k(w)=ik\,c_k(u)$, et $\sum_kk^2|c_k(u)|^2=\norm w_H^2<\infty$ : $u\in H^1_\sharp$. Ainsi
\[
H^1_\sharp=\{u\in L^2(J):\ (u^\sharp)'\in L^2_{\mathrm{loc}}(\R)\},
\]
et $u'$ est bien « la » dérivée de $u$, au sens des distributions \emph{périodiques}.

\medskip
\noindent\textbf{(c)} $c_0(f)=\frac1{2\pi}\int_0^{2\pi}x\,dx=\pi$ et, pour $k\ne0$, par parties,
\[
c_k(f)=\frac1{2\pi}\int_0^{2\pi}xe^{-ikx}\,dx=\frac1{2\pi}\Big(\Big[\frac{xe^{-ikx}}{-ik}\Big]_0^{2\pi}+\frac1{ik}\int_0^{2\pi}e^{-ikx}\,dx\Big)=\frac1{2\pi}\cdot\frac{2\pi}{-ik}=\frac ik .
\]
Donc $k^2|c_k(f)|^2=1$ et $\sum_kk^2|c_k(f)|^2=\infty$ : $f\notin H^1_\sharp$. Pourtant $f$ est lisse sur $J$, de dérivée classique $1$, donc $T_f'=T_1$ dans $\D'(J)$ (Proposition~6.3).

\emph{Calcul de $(f^\sharp)'$.} \emph{Première façon :} $f^\sharp$ est la fonction « dents de scie », $f^\sharp(x)=x-2\pi n$ sur $]2\pi n,2\pi(n+1)[$ : elle est $\mathcal C^1$ par morceaux, de dérivée $1$, avec un saut $f^\sharp((2\pi n)^+)-f^\sharp((2\pi n)^-)=0-2\pi=-2\pi$ en chaque point de $2\pi\Z$. La formule des sauts (version localement finie, Exercice~4(a)) donne $(f^\sharp)'=1-2\pi\Sha$, avec $\Sha=\sum_n\delta_{2\pi n}$. \emph{Seconde façon :} par (a), $(f^\sharp)'=\lim_N\sum_{0<|k|\le N}ik\cdot\frac ik\,e_k=\lim_N\big(1-\sum_{|k|\le N}e_k\big)=1-2\pi\Sha$ par le TD0 ch.~6 (peigne de Dirac). Les deux calculs concordent.

\emph{Commentaire.} La dérivée de $f$ dans $\D'(J)$ est dans $L^2$, mais celle de $f^\sharp$ dans $\D'(\R)$ ne l'est pas : le recollement périodique crée un saut, donc une masse de Dirac, invisible depuis l'ouvert $J$ (aucune fonction test de $\D(J)$ ne voit les points $0$ et $2\pi$). La condition $\sum k^2|c_k|^2<\infty$ caractérise les fonctions de $L^2(J)$ à dérivée $L^2$ \emph{et} « sans saut entre $2\pi^-$ et $0^+$ », ce que (b) exprime par $\tilde u(0)=\tilde u(2\pi)$.

\medskip
\noindent\textbf{(d)} \emph{Existence.} Posons $\varphi_f:=-\sum_{k\ne0}k^{-2}c_k(f)e_k$ : la série converge dans $H$ (Proposition~4.4, $\sum k^{-4}|c_k(f)|^2\le\norm f_H^2$), et $c_0(\varphi_f)=0$, donc $\varphi_f\in L^2_0$. Par (a), $S_N\varphi_f\wk\varphi_f^\sharp$ dans $\D'(\R)$, donc $(S_N\varphi_f)''\wk(\varphi_f^\sharp)''$. Or
\[
(S_N\varphi_f)''=-\sum_{0<|k|\le N}\frac{(ik)^2}{k^2}c_k(f)e_k=\sum_{0<|k|\le N}c_k(f)e_k=S_Nf
\]
(car $c_0(f)=0$), et $S_Nf\wk f^\sharp$ par (a). Donc $(\varphi_f^\sharp)''=f^\sharp$.

\emph{Unicité.} Soient $\varphi_1,\varphi_2\in L^2_0$ deux solutions et $w:=\varphi_1-\varphi_2\in L^2_0$ : $(w^\sharp)''=0$ dans $\D'(\R)$. Par le Théorème~6.4 ($\R$ connexe), $(w^\sharp)'=c_1$ constante ; puis $(w^\sharp-c_1x)'=0$, donc $w^\sharp=c_1x+c_0$ presque partout. La périodicité de $w^\sharp$ impose $c_1x+c_0=c_1(x+2\pi)+c_0$ p.p., soit $c_1=0$ ; puis $c_0=c_0(w)=0$. Donc $w=0$.

\emph{Régularité.} $c_k(\varphi_f)=-c_k(f)/k^2$, donc $\sum k^2|c_k(\varphi_f)|^2<\infty$ : $\varphi_f\in H^1_\sharp$, avec $c_k(\varphi_f')=ik\,c_k(\varphi_f)=-ic_k(f)/k$, donc $\sum k^2|c_k(\varphi_f')|^2=\norm f^2_H<\infty$ : $\varphi_f'\in H^1_\sharp$ aussi. Par (b), les séries de Fourier de $\varphi_f$ et de $\varphi_f'$ convergent normalement : la série de $\varphi_f$ converge uniformément ainsi que sa série dérivée, donc sa somme est de classe $\mathcal C^1$ (théorème de dérivation terme à terme), $2\pi$-périodique.

\emph{Non-unicité dans $\D'(J)$.} La fonction $g(x):=x-\pi$ est dans $L^2_0$ ($c_0(g)=\pi-\pi=0$) et lisse sur $J$ avec $g''=0$. Pour tout $a\in\C$, $\varphi_f+ag\in L^2_0$ et $(\varphi_f+ag)''=f$ dans $\D'(J)$ : il y a une droite de solutions. (Elles ne sont pas périodiques : $(g^\sharp)''=-2\pi\Sha'\ne0$ par (c).)

\medskip
\noindent\textbf{(e)} $L^2_0=\{e_0\}^\perp$ est un sous-espace fermé de $H$, donc un espace de Hilbert, de base hilbertienne $(e_k)_{k\in\Z^\star}$. La formule $Kf=-\sum_{k\ne0}k^{-2}c_k(f)e_k$ montre que $K$ est linéaire, et par Parseval
\[
\norm{Kf}_H^2=\sum_{k\ne0}\frac{|c_k(f)|^2}{k^4}\le\sum_{k\ne0}|c_k(f)|^2=\norm f_H^2 ,
\]
avec égalité pour $f=e_1$ : $K$ est continu et $\opnorm K=1$.

\emph{Auto-adjoint.} Par Parseval polarisé, $\langle Kf,g\rangle=-\sum_{k\ne0}k^{-2}c_k(f)\overline{c_k(g)}=\langle f,Kg\rangle$ (les coefficients $-k^{-2}$ sont réels). Donc $K^\star=K$.

\emph{Compact.} Soit $K_mf:=-\sum_{0<|k|\le m}k^{-2}c_k(f)e_k$, opérateur continu de rang fini. Par Parseval, $\norm{(K-K_m)f}^2=\sum_{|k|>m}k^{-4}|c_k(f)|^2\le(m+1)^{-4}\norm f^2$, donc $\opnorm{K-K_m}\le(m+1)^{-2}$. Soit $B$ la boule unité fermée de $L^2_0$ et $\varepsilon>0$ ; choisissons $m$ avec $(m+1)^{-2}\le\varepsilon$. L'ensemble $K_m(B)$ est une partie bornée de l'espace de dimension finie $\Vect(e_k:0<|k|\le m)$, donc précompacte : elle est recouverte par un nombre fini de boules $B(y_j,\varepsilon)$. Pour $f\in B$, $\norm{Kf-K_mf}\le\varepsilon$, donc $Kf\in B(y_j,2\varepsilon)$ pour un $j$ : $K(B)$ est recouvert par un nombre fini de boules de rayon $2\varepsilon$. Ainsi $K(B)$ est précompacte dans l'espace complet $L^2_0$, donc d'adhérence compacte : $K$ est compact (Remarque~4.7).

\emph{Éléments propres.} $Ke_k=-k^{-2}e_k$. Si $Kf=\lambda f$, alors $(-k^{-2}-\lambda)c_k(f)=0$ pour tout $k\ne0$ : ou bien $\lambda=-1/n^2$ pour un $n\ge1$, et alors $c_k(f)=0$ sauf pour $k=\pm n$, ou bien $f=0$. Les valeurs propres sont donc les $-1/n^2$, $n\ge1$, d'espace propre $\Vect(e_n,e_{-n})$, de dimension~$2$ (engendré aussi par $\cos(nx)$ et $\sin(nx)$) ; $0$ n'est pas valeur propre ($K$ est injectif). C'est conforme à la Proposition~4.7 : valeurs propres réelles, espaces propres orthogonaux et de dimension finie, $0$ seul point d'accumulation ; et au Théorème~5.5 : $(e_k)_{k\ne0}$ est une base hilbertienne de vecteurs propres. Noter que $\sup|\lambda|=1=\opnorm K$.

\medskip
\noindent\textbf{(f)} \emph{Structure hilbertienne.} Par Parseval polarisé et $c_k(u')=ik\,c_k(u)$,
\[
\langle u,v\rangle_1=\sum_k(1+k^2)c_k(u)\overline{c_k(v)} .
\]
C'est un produit scalaire sur $H^1_\sharp$, et l'application
\[
\Lambda:H^1_\sharp\to\ell^2(\Z),\qquad u\mapsto\big(\sqrt{1+k^2}\,c_k(u)\big)_{k\in\Z}
\]
est linéaire et isométrique. Elle est surjective : si $(a_k)\in\ell^2(\Z)$, la série $u:=\sum_k(1+k^2)^{-1/2}a_ke_k$ converge dans $H$ (Proposition~4.4), $c_k(u)=(1+k^2)^{-1/2}a_k$ et $\sum k^2|c_k(u)|^2\le\sum|a_k|^2<\infty$, donc $u\in H^1_\sharp$ et $\Lambda u=(a_k)$. Ainsi $(H^1_\sharp,\langle\cdot,\cdot\rangle_1)$ est isométriquement isomorphe à $\ell^2(\Z)$ : c'est un espace de Hilbert séparable.

\emph{Riesz.} Soit $f\in H$. La forme $L:v\mapsto\langle v,f\rangle$ est linéaire sur $H^1_\sharp$ et continue : $|L(v)|\le\norm f_H\norm v_H\le\norm f_H\norm v_1$. Par le Théorème~4.6, il existe un unique $u\in H^1_\sharp$ tel que $L(v)=\langle v,u\rangle_1$ pour tout $v\in H^1_\sharp$. Avec $v=e_k\in H^1_\sharp$ : $\overline{c_k(f)}=\langle e_k,f\rangle=\langle e_k,u\rangle_1=(1+k^2)\overline{c_k(u)}$, d'où
\[
c_k(u)=\frac{c_k(f)}{1+k^2},\qquad u=\sum_{k\in\Z}\frac{c_k(f)}{1+k^2}\,e_k .
\]
\emph{Équation.} Par (a), $S_Nu\wk u^\sharp$ dans $\D'(\R)$, donc $-(S_Nu)''+S_Nu\wk-(u^\sharp)''+u^\sharp$. Or $-(S_Nu)''+S_Nu=\sum_{|k|\le N}(k^2+1)\frac{c_k(f)}{1+k^2}e_k=S_Nf\wk f^\sharp$. Donc $-(u^\sharp)''+u^\sharp=f^\sharp$ dans $\D'(\R)$.
\end{solution}

\begin{piege}
\textbf{Dans quel espace de distributions lit-on l'équation ?} La Proposition~6.6 parle de l'équation $\varphi''=f$ « au sens des distributions » et affirme l'unicité dans $L^2_0(J_{2\pi})$. La question (d) montre que cette unicité n'est vraie que pour les prolongements périodiques, dans $\D'(\R)$ : dans $\D'(J)$, on peut toujours ajouter $a(x-\pi)$. La raison est celle de (c) : une série trigonométrique peut converger vers $0$ dans $\D'(J)$ sans que ses coefficients soient nuls ($\sum_{k}e_k\wk2\pi\Sha$, qui est nulle sur $J$), de sorte que « l'unicité de la décomposition » invoquée dans la preuve du polycopié n'a de sens que dans $L^2$ ou pour les distributions périodiques. À l'examen, préciser l'ouvert ($J$ ou $\R$) chaque fois que l'on dérive une série de Fourier.
\end{piege}

\begin{pointcle}
\textbf{Le dictionnaire de Fourier, d'un chapitre à l'autre.} Sur les coefficients $(c_k)$ : appartenir à $L^2$, c'est $(c_k)\in\ell^2$ (chapitre~4, Parseval) ; dériver, c'est multiplier par $ik$ (chapitre~6, continuité de la dérivation dans $\D'$) ; résoudre $\varphi''=f$, c'est diviser par $-k^2$, ce qui donne un opérateur diagonal de valeurs propres tendant vers $0$, donc compact et auto-adjoint (chapitres~4 et 5) ; et l'existence d'une solution faible de $-u''+u=f$ est le théorème de représentation de Riesz dans l'espace de Hilbert $H^1_\sharp$. Un même calcul formel sur les $c_k$ est justifié trois fois, par trois théorèmes différents : c'est ce type de question « transversale » qui clôt naturellement un sujet d'examen.
\end{pointcle}


% =====================================================================
\begin{exo}{Exercice 13 -- Différences secondes, convexité, équation des ondes}
\textbf{A. Convexité (Exemple~6.11).} Pour $T\in\D'(\R)$ et $h\ne0$ on pose $D_hT:=\dfrac{\tau_hT+\tau_{-h}T-2T}{h^2}$.
\begin{enumerate}[label=(\alph*)]
  \item Montrer que $D_hT\wk T''$ dans $\D'(\R)$ quand $h\to0$. En déduire la limite de $n^2\big(\delta_{1/n}+\delta_{-1/n}-2\delta_0\big)$.
  \item Soit $f:\R\to\R$ convexe (donc continue). Montrer que $T_f''$ est une distribution positive, donc une mesure de Radon positive.
  \item Réciproquement, soit $f:\R\to\R$ continue telle que $T_f''$ soit positive. Montrer que $f$ est convexe. \emph{Indication : régulariser par $f\star\rho_n$.}
  \item Les fonctions $|x|$, $e^{-|x|}$ et $|x^2-1|$ sont-elles convexes ? Répondre en calculant leur dérivée seconde dans $\D'(\R)$.
\end{enumerate}
\textbf{B. Équation des ondes.} On fixe $c>0$ et l'on note $\square:=\partial_{tt}-c^2\partial_{xx}$ sur $\R^2$, de variables $(x,t)$.
\begin{enumerate}[label=(\alph*),resume]
  \item Soit $f\in L^1_{\mathrm{loc}}(\R)$ et $u(x,t):=f(x-ct)$. Montrer que $u\in L^1_{\mathrm{loc}}(\R^2)$, que $\partial_tu+c\,\partial_xu=0$ dans $\D'(\R^2)$, puis que $\square u=0$. L'opérateur $\square$ vérifie-t-il un analogue du lemme de Weyl ?
  \item Soit $E:=\frac1{2c}\mathbf 1_{\mathcal C}$ avec $\mathcal C:=\{(x,t):t>0,\ |x|<ct\}$. Montrer que $\square E=\delta_{(0,0)}$ dans $\D'(\R^2)$.
  \item Soit $\psi\in\D(\R^2)$. Expliciter $u:=\psi\star E$, montrer que $\square u=\psi$ et que $u(x,t)=0$ si $t\le\min\{s:\ (y,s)\in\supp\psi\}$.
\end{enumerate}
\end{exo}

\begin{solution}
\noindent\textbf{(a)} Pour $\varphi\in\D(\R)$, la Définition~6.7 donne $\langle\tau_{\pm h}T,\varphi\rangle=\langle T,\tau_{\mp h}\varphi\rangle$, donc $\langle D_hT,\varphi\rangle=\langle T,D_h\varphi\rangle$ avec
\[
D_h\varphi(x):=\frac{\varphi(x+h)+\varphi(x-h)-2\varphi(x)}{h^2}.
\]
Montrons que $D_h\varphi\to\varphi''$ dans $\D(\R)$ quand $h\to0$. La formule de Taylor avec reste intégral, $\varphi(x\pm h)=\varphi(x)\pm h\varphi'(x)+h^2\int_0^1(1-s)\varphi''(x\pm sh)\,ds$, donne, puisque $\int_0^1(1-s)\,ds=\frac12$,
\[
D_h\varphi(x)-\varphi''(x)=\int_0^1(1-s)\big(\varphi''(x+sh)+\varphi''(x-sh)-2\varphi''(x)\big)\,ds ,
\]
majoré en module par $\int_0^1(1-s)\,2s|h|\norm{\varphi'''}_\infty\,ds\le|h|\norm{\varphi'''}_\infty$ (accroissements finis). Comme $(D_h\varphi)^{(j)}=D_h(\varphi^{(j)})$, le même calcul appliqué à $\varphi^{(j)}$ donne, pour tout $j$,
\[
\norm{(D_h\varphi)^{(j)}-\varphi^{(j+2)}}_\infty\le|h|\,\norm{\varphi^{(j+3)}}_\infty\xrightarrow[h\to0]{}0 .
\]
Enfin, pour $|h|\le1$, $\supp D_h\varphi\subset\supp\varphi+[-1,1]$, compact fixe. Donc, pour toute suite $h_n\to0$, $D_{h_n}\varphi\to\varphi''$ dans $\D(\R)$, et la continuité séquentielle de $T$ donne $\langle D_{h_n}T,\varphi\rangle\to\langle T,\varphi''\rangle=\langle T'',\varphi\rangle$ : $D_hT\wk T''$. (C'est la version à l'ordre~$2$ des Propositions~6.1 et 6.7.)

Pour $T=\delta_0$ : $\tau_a\delta_0=\delta_a$, donc $n^2(\delta_{1/n}+\delta_{-1/n}-2\delta_0)=D_{1/n}\delta_0\wk\delta_0''$. Directement : $n^2\big(\varphi(\tfrac1n)+\varphi(-\tfrac1n)-2\varphi(0)\big)\to\varphi''(0)=\langle\delta_0'',\varphi\rangle$.

\medskip
\noindent\textbf{(b)} Par l'Exercice~6.3, $\tau_{\pm h}T_f=T_{\tau_{\pm h}f}$, donc $D_hT_f=T_{D_hf}$ avec $D_hf(x)=\frac{f(x+h)+f(x-h)-2f(x)}{h^2}$. La convexité de $f$ donne $f(x)=f\big(\frac{(x+h)+(x-h)}2\big)\le\frac{f(x+h)+f(x-h)}2$, \ie\ $D_hf\ge0$ sur $\R$. Pour $\varphi\in\D(\R)$ positive, $\langle D_hT_f,\varphi\rangle=\int D_hf\,\varphi\ge0$, et en passant à la limite avec (a) : $\langle T_f'',\varphi\rangle\ge0$. Donc $T_f''$ est une forme linéaire positive sur $\D(\R)$ : par le Théorème~6.2 c'est une distribution d'ordre $0$, une mesure de Radon positive (Définition~6.5).

\medskip
\noindent\textbf{(c)} Soit $(\rho_n)$ une approximation de l'unité issue d'une fonction test $\rho\ge0$, et $f_n:=f\star\rho_n\in\Cinf(\R)$. On a $T_{f_n}=\rho_n\star T_f$ (Exercice~6.4) et, par la Proposition~6.10, $T_{f_n''}=(\rho_n\star T_f)''=\rho_n\star T_f''$. Pour $\varphi\in\D(\R)$ positive, $\varphi\star\check\rho_n$ est une fonction test positive, donc
\[
\int_\R f_n''\varphi=\langle\rho_n\star T_f'',\varphi\rangle=\langle T_f'',\varphi\star\check\rho_n\rangle\ge0 .
\]
La fonction continue $f_n''$ est donc positive : sinon $f_n''<0$ sur un intervalle ouvert, et une bosse $\varphi\ge0$ non nulle supportée dans cet intervalle donnerait $\int f_n''\varphi<0$. Ainsi $f_n$ est convexe (fonction $\mathcal C^2$ à dérivée seconde positive). Enfin $f_n\to f$ simplement (et même uniformément sur tout compact, $f$ étant continue : Théorème~3.6 appliqué à $f\theta$, $\theta$ plateau), et une limite simple de fonctions convexes est convexe (passer à la limite dans $f_n(\lambda x+(1-\lambda)y)\le\lambda f_n(x)+(1-\lambda)f_n(y)$). Donc $f$ est convexe.

\medskip
\noindent\textbf{(d)} $|x|''=2\delta_0$ (TD0 ch.~6) : positive, donc $|x|$ est convexe. $(e^{-|x|})''=e^{-|x|}-2\delta_0$ (Exercice~4(a)) : pour $\varphi\ge0$ concentrée près de $0$, par exemple $\varphi=\chi(n\,\cdot)$ avec $\chi\ge0$, $\chi(0)=1$, on a $\int e^{-|x|}\varphi\le\frac1n\norm\chi_1<2=2\varphi(0)$ pour $n$ grand : la distribution n'est pas positive, $e^{-|x|}$ n'est pas convexe (elle a un « pic » en $0$). Pour $g(x)=|x^2-1|$ : continue, $g'=2x\,\sgn(x^2-1)$ hors de $\pm1$, de sauts $g'(1^+)-g'(1^-)=2-(-2)=4$ et $g'((-1)^+)-g'((-1)^-)=2-(-2)=4$, et de dérivée $2\sgn(x^2-1)$ ailleurs. Donc $g''=2\sgn(x^2-1)+4\delta_1+4\delta_{-1}$, qui vaut $-2$ sur $]-1,1[$ : pour $\varphi\ge0$ non nulle à support dans $]-1,1[$, $\langle g'',\varphi\rangle=-2\int\varphi<0$. $g$ n'est pas convexe (elle est concave sur $[-1,1]$).

\medskip
\noindent\textbf{(e)} \emph{Intégrabilité locale.} Pour $R>0$, par invariance par translation de l'intégrale (Proposition~A.4), $\int_{-R}^R\int_{-R}^R|f(x-ct)|\,dx\,dt\le\int_{-R}^R\int_{|y|\le(1+c)R}|f(y)|\,dy\,dt=2R\norm f_{L^1([-(1+c)R,(1+c)R])}<\infty$.

\emph{Équation de transport.} Soit $\varphi\in\D(\R^2)$ à support dans $[-R,R]^2$. Par définition, puis par le changement de variable $x=y+ct$ à $t$ fixé et le théorème de Fubini (l'intégrande est dominé par $C|f(y)|\mathbf 1_{|y|\le(1+c)R}\mathbf 1_{|t|\le R}$, intégrable),
\begin{align*}
\langle\partial_tu+c\,\partial_xu,\varphi\rangle&=-\int_\R\int_\R f(x-ct)(\partial_t\varphi+c\,\partial_x\varphi)(x,t)\,dx\,dt\\
&=-\int_\R f(y)\Big(\int_\R(\partial_t\varphi+c\,\partial_x\varphi)(y+ct,t)\,dt\Big)dy .
\end{align*}
Or $(\partial_t\varphi+c\,\partial_x\varphi)(y+ct,t)=\frac d{dt}\big[\varphi(y+ct,t)\big]$, et $t\mapsto\varphi(y+ct,t)$ est lisse à support compact : l'intégrale intérieure est nulle. Donc $\partial_tu+c\,\partial_xu=0$.

\emph{Équation des ondes.} Les dérivées partielles des distributions commutent (lemme de Schwarz pour les fonctions tests), donc $\square=(\partial_t-c\,\partial_x)(\partial_t+c\,\partial_x)$ sur $\D'(\R^2)$, et $\square u=(\partial_t-c\,\partial_x)0=0$. De même $f(x+ct)$ est solution.

\emph{Pas de lemme de Weyl.} Avec $f=h$, $u=\mathbf 1_{x\ge ct}$ est une solution de $\square u=0$ qui n'est pas une fonction lisse, ni même continue ; avec $f\in L^1_{\mathrm{loc}}$ quelconque on obtient des solutions arbitrairement irrégulières. Les singularités de la donnée se propagent le long des droites $x\pm ct=\text{cte}$ au lieu d'être effacées : $\square$ n'est pas hypoelliptique, contrairement à $\Delta$ (Lemme~6.4, Remarque~6.9).

\medskip
\noindent\textbf{(f)} $E$ est bornée, donc dans $L^1_{\mathrm{loc}}(\R^2)$. Soit $\varphi\in\D(\R^2)$. Deux dérivations dans chaque terme : $\langle\square E,\varphi\rangle=\langle E,\square\varphi\rangle=\frac1{2c}\iint_{\mathcal C}(\partial_{tt}\varphi-c^2\partial_{xx}\varphi)$. On calcule chaque morceau par Fubini, en intégrant d'abord dans la direction de la dérivée. À $x$ fixé, la tranche de $\mathcal C$ est $\{t>|x|/c\}$ :
\begin{align*}
\iint_{\mathcal C}\partial_{tt}\varphi&=\int_\R\Big(\int_{|x|/c}^{+\infty}\partial_{tt}\varphi(x,t)\,dt\Big)dx=-\int_\R\partial_t\varphi\big(x,\tfrac{|x|}c\big)\,dx\\
&=-c\int_0^{+\infty}\big(\partial_t\varphi(cs,s)+\partial_t\varphi(-cs,s)\big)\,ds ,
\end{align*}
en posant $x=cs$ sur $\R_+$ et $x=-cs$ sur $\R_-$. À $t>0$ fixé, la tranche est $\{|x|<ct\}$ :
\[
\iint_{\mathcal C}\partial_{xx}\varphi=\int_0^{+\infty}\Big(\int_{-ct}^{ct}\partial_{xx}\varphi(x,t)\,dx\Big)dt=\int_0^{+\infty}\big(\partial_x\varphi(cs,s)-\partial_x\varphi(-cs,s)\big)\,ds .
\]
Donc
\begin{align*}
\langle\square E,\varphi\rangle&=-\frac12\int_0^{+\infty}\Big[(\partial_t\varphi+c\,\partial_x\varphi)(cs,s)+(\partial_t\varphi-c\,\partial_x\varphi)(-cs,s)\Big]ds\\
&=-\frac12\int_0^{+\infty}\frac d{ds}\big[\varphi(cs,s)+\varphi(-cs,s)\big]\,ds ,
\end{align*}
ce qui vaut $-\frac12\big(0-2\varphi(0,0)\big)=\varphi(0,0)$, $\varphi$ étant à support compact. Donc $\square E=\delta_{(0,0)}$ : $E$ est une solution fondamentale de l'opérateur des ondes, portée par le « cône de lumière futur » $\mathcal C$.

\medskip
\noindent\textbf{(g)} Par l'Exercice~6.4, $\psi\star T_E=T_{\psi\star E}$, avec
\[
u(x,t)=(\psi\star E)(x,t)=\iint\psi(x-y,t-s)E(y,s)\,dy\,ds=\frac1{2c}\iint_{\{s<t,\ |x-y|<c(t-s)\}}\psi(y,s)\,dy\,ds
\]
(changement de variable $(y,s)\mapsto(x-y,t-s)$). Par la Proposition~6.10 (en dimension~$2$), $u$ est lisse et $\square u=\psi\star\square E=\psi\star\delta_{(0,0)}=\psi$ (Exemple~6.13). Si $t\le t_0:=\min\{s:(y,s)\in\supp\psi\}$, le domaine d'intégration $\{s<t\}$ ne rencontre pas $\supp\psi$ : $u(x,t)=0$. C'est la \emph{causalité} : la solution $\psi\star E$ est nulle avant que la source ne s'allume ; et la valeur de $u$ en $(x,t)$ ne dépend que de la source dans le cône passé $\{|x-y|<c(t-s)\}$ : l'information se propage à vitesse au plus $c$.
\end{solution}

\begin{pointcle}
\textbf{Positivité et régularité se lisent sur les dérivées distributionnelles.} « $T'=0\iff T$ constante » (Théorème~6.4), « $T'\ge0\iff T$ croissante », « $T''\ge0\iff T$ convexe » (Exemple~6.11) : les caractérisations de l'analyse réelle survivent, à condition d'autoriser des mesures comme dérivées ($|x|''=2\delta_0$). Le schéma de preuve est toujours le même : dans un sens, un quotient différentiel positif qui converge dans $\D'$ (question (a)) ; dans l'autre, une régularisation $f\star\rho_n$ qui préserve la positivité parce que $\rho_n\ge0$.
\end{pointcle}

\begin{commentaire}
\textbf{Elliptique contre hyperbolique.} Le laplacien et l'opérateur des ondes ne diffèrent que par un signe, mais leurs solutions fondamentales n'ont rien de commun : $-\frac1{4\pi\norm x}$ est lisse hors de l'origine, d'où le lemme de Weyl ; $\frac1{2c}\mathbf 1_{\mathcal C}$ est discontinue le long de tout le cône $|x|=ct$, et les solutions de $\square u=0$ peuvent être aussi irrégulières que l'on veut (question (e)). En contrepartie, le support de $E$ donne la causalité et la vitesse finie de propagation, exactement comme $h\cdot y$ en dimension~$1$ (Exercice~8). Le calcul de (f) n'utilise que Fubini et l'intégration d'une dérivée : c'est la version bidimensionnelle de la formule des sauts.
\end{commentaire}


% =====================================================================
\begin{exo}{Exercice 14 -- Vrai ou faux ?}
Pour chaque affirmation, dire si elle est vraie ou fausse, et justifier par une preuve courte ou un contre-exemple. (Viser deux à trois minutes par question.)
\begin{enumerate}[label=\arabic*.]
  \item Si $(\varphi_n)$ est une suite de $\D(\R)$ telle que $\varphi_n^{(j)}\to0$ uniformément sur $\R$ pour tout $j\in\N$, alors $\varphi_n\to0$ dans $\D(\R)$.
  \item Si $T\in\D'(\R)$ et si $\varphi\in\D(\R)$ est nulle sur $\supp T$, alors $\langle T,\varphi\rangle=0$.
  \item Si $(f_n)$ est une suite de $L^1_{\mathrm{loc}}(\R)$ qui converge presque partout vers $f\in L^1_{\mathrm{loc}}(\R)$, alors $T_{f_n}\wk T_f$.
  \item Si $T_{f_n}\wk0$ dans $\D'(\R)$, alors une sous-suite de $(f_n)$ converge vers $0$ presque partout.
  \item $\delta_0'$ est une mesure de Radon.
  \item Si $T\in\D'(\R)$ vérifie $xT=0$, alors $T=0$.
  \item Si $f\in\Cinf(\R)$ ne s'annule pas et si $fT=0$, alors $T=0$.
  \item $(x\delta_0')'=-\delta_0'$.
  \item On peut prolonger le produit $(f,T)\mapsto fT$ en donnant un sens à $h\,\delta_0$ par continuité : si $T_{g_n}\wk\delta_0$ avec $g_n\in\D(\R)$, la suite $T_{hg_n}$ a une limite indépendante du choix de $(g_n)$.
  \item Les solutions dans $\D'(\R)$ de $T'=T$ sont les $Ce^x$, $C\in\C$.
  \item Si $\psi\in\D(\R)$ est non nulle et $\psi\star T=0$, alors $T=0$.
  \item Si $f\in L^1_{\mathrm{loc}}(\R)$ et si $\langle T_f',\varphi\rangle=0$ pour toute $\varphi\in\D(\R\setminus\{0\})$, alors $f$ est constante presque partout.
  \item Si $T_n'\wk0$ dans $\D'(\R)$, alors $(T_n)$ converge dans $\D'(\R)$.
  \item La dérivée de $\ln|x|$ dans $\D'(\R)$ est la fonction $x\mapsto1/x$.
  \item Si $T\in\D'(\R)$ est d'ordre $m$, alors $T'$ est d'ordre exactement $m+1$.
  \item Si $T\in\D'(\R)$ et si $T'$ est d'ordre $0$, alors $T$ est une fonction, et même $T\in L^\infty_{\mathrm{loc}}(\R)$.
\end{enumerate}
\end{exo}

\begin{solution}
\noindent\textbf{1. Faux.} Il manque la condition de support (Définition~6.2). Soit $\varphi\in\D(]-\frac12,\frac12[)$ avec $\varphi(0)=1$ et $\varphi_n:=\frac1n\varphi(\cdot-n)$ : $\norm{\varphi_n^{(j)}}_\infty=\frac1n\norm{\varphi^{(j)}}_\infty\to0$ pour tout $j$, mais les supports partent à l'infini. La distribution $T:=\sum_{k\ge1}k\,\delta_k$ (somme localement finie, ordre $0$) vérifie $T(\varphi_n)=n\cdot\frac1n\varphi(0)=1\not\to0=T(0)$ : la suite ne converge pas dans $\D(\R)$, sans quoi la continuité séquentielle de $T$ serait violée.

\medskip
\noindent\textbf{2. Faux.} $T=\delta_0'$, de support $\{0\}$, et $\varphi(x)=x\chi(x)$ avec $\chi\in\D(\R)$ valant $1$ près de $0$ : $\varphi(0)=0$ mais $\langle\delta_0',\varphi\rangle=-\varphi'(0)=-1$. L'énoncé correct demande $\varphi$ nulle \emph{au voisinage} de $\supp T$, \ie\ $\supp\varphi\cap\supp T=\emptyset$ (Exercice~6(c)).

\medskip
\noindent\textbf{3. Faux.} $f_n=n\mathbf 1_{]0,1/n[}\to0$ en tout point, mais $T_{f_n}\wk\delta_0\ne0$ (Exercice~3(a)). Il manque une domination : la convergence p.p.\ avec $|f_n|\le g\in L^1_{\mathrm{loc}}$ suffirait (convergence dominée).

\medskip
\noindent\textbf{4. Faux.} $f_n(x)=e^{inx}$ : $T_{f_n}\wk0$ par Riemann--Lebesgue, mais $|f_n|=1$ partout. (Version réelle : $\sin(nx)\wk0$ ; si $\sin(n_kx)\to0$ p.p., alors par convergence dominée $\int_0^{2\pi}\sin^2(n_kx)\,dx\to0$, alors que cette intégrale vaut $\pi$.)

\medskip
\noindent\textbf{5. Faux.} Une mesure de Radon est d'ordre $0$ (Théorème~6.2), et $\delta_0'$ est d'ordre exactement $1$ (TD0 ch.~6, Exemple~6.4(b)).

\medskip
\noindent\textbf{6. Faux.} $x\delta_0=0$ et $\delta_0\ne0$. Les solutions sont exactement les $c\,\delta_0$ (TD0 ch.~6).

\medskip
\noindent\textbf{7. Vrai.} $1/f\in\Cinf(\R)$, et par associativité du produit par les fonctions lisses, $T=\big(\frac1f\cdot f\big)T=\frac1f(fT)=0$ : $\langle T,\varphi\rangle=\langle T,f\cdot\frac\varphi f\rangle=\langle fT,\frac\varphi f\rangle=0$.

\medskip
\noindent\textbf{8. Vrai.} $x\delta_0'=-\delta_0$ (TD0 ch.~6), donc $(x\delta_0')'=-\delta_0'$. Par Leibniz : $(x\delta_0')'=\delta_0'+x\delta_0''=\delta_0'-2\delta_0'$, car $x\delta_0''=-2\delta_0'$ (Exercice~6(a)).

\medskip
\noindent\textbf{9. Faux.} Soit $\rho\in\D(\R)$ positive d'intégrale $1$ et $\rho_n=n\rho(n\,\cdot)\wk\delta_0$. Si $\rho$ est paire, $\int h\rho_n\varphi=\int_0^\infty\rho(u)\varphi(u/n)\,du\to\varphi(0)\int_0^\infty\rho=\frac12\varphi(0)$ : $h\rho_n\wk\frac12\delta_0$. Si $\rho$ est à support dans $]0,1[$, $h\rho_n=\rho_n\wk\delta_0$ ; si elle est à support dans $]-1,0[$, $h\rho_n=0$. La limite dépend de la suite : $h\,\delta_0$ n'a pas de sens. Le produit n'est défini que si l'un des facteurs est lisse (Proposition~6.8).

\medskip
\noindent\textbf{10. Vrai.} Par Leibniz, $(e^{-x}T)'=e^{-x}(T'-T)=0$, donc $e^{-x}T=C$ (Théorème~6.4) et $T=Ce^x$ ; réciproquement $(Ce^x)'=Ce^x$. Pas de solution « exotique » : l'équation n'a pas de point singulier, contrairement à $xT'+T=0$ (Exercice~7(c)).

\medskip
\noindent\textbf{11. Faux.} $T=T_1$ (constante $1$) : $\langle\psi\star1,\varphi\rangle=\int(\varphi\star\check\psi)=\big(\int\varphi\big)\big(\int\psi\big)$, donc $\psi\star1$ est la constante $\int\psi$, nulle pour $\psi=\rho'$ avec $\rho\in\D(\R)$ non nulle.

\medskip
\noindent\textbf{12. Faux.} $f=h$ : $T_h'=\delta_0$ est nulle sur $\R\setminus\{0\}$, et $h$ n'est pas constante. Le Théorème~6.4 demande un ouvert connexe : on obtient seulement $f$ constante sur $]-\infty,0[$ et sur $]0,+\infty[$.

\medskip
\noindent\textbf{13. Faux.} $T_n:=(-1)^n$ (constantes) : $T_n'=0$, et $\langle T_n,\varphi\rangle=(-1)^n\int\varphi$ diverge si $\int\varphi\ne0$. La dérivation perd les constantes.

\medskip
\noindent\textbf{14. Faux.} $x\mapsto1/x$ n'est pas localement intégrable près de $0$, elle ne définit pas une distribution sur $\R$. La dérivée est $\mathrm{vp}(1/x)$ (Proposition~6.4), qui n'est ni une fonction ni une mesure (ordre $1$).

\medskip
\noindent\textbf{15. Faux.} L'ordre augmente \emph{d'au plus} $1$ (Proposition~6.3). $T=|x|$ est d'ordre $0$ et $T'=\sgn$ aussi ; de même $h'=\delta_0$.

\medskip
\noindent\textbf{16. Vrai.} Notons $\mu:=T'$. Par l'Exercice~5, $T=S+c$ avec $\langle S,\varphi\rangle=-\langle\mu,\Psi\rangle$, où $\Psi$ est la primitive à support compact de $\psi=\varphi-(\int\varphi)\theta$, à support dans $\tilde K=\mathrm{conv}(K\cup\supp\theta)$ si $\varphi\in\D_K(\R)$. Comme $\mu$ est d'ordre $0$,
\[
|\langle S,\varphi\rangle|\le C_{\tilde K}\norm\Psi_\infty\le C_{\tilde K}\norm\psi_1\le C_{\tilde K}\big(1+\norm\theta_1\big)\norm\varphi_1\qquad(\varphi\in\D_K(\R)),
\]
car $|\Psi(x)|=|\int_{-\infty}^x\psi|\le\norm\psi_1$ et $\norm\psi_1\le\norm\varphi_1+|\int\varphi|\norm\theta_1$. C'est l'hypothèse de l'Exercice~6.2 du polycopié avec $p=1$ : $S\in L^\infty_{\mathrm{loc}}(\R)$, donc $T$ aussi. Exemple : $\delta_0$ a pour primitives $h+c$. (Ce sont les fonctions à variation localement bornée.)
\end{solution}

\begin{piege}
\textbf{Les six erreurs que ces questions cherchent.} (i) Oublier la condition de support dans la convergence de $\D$ (question~1). (ii) Confondre « nulle sur le support » et « nulle au voisinage du support » (2). (iii) Croire que la convergence dans $\D'$ est comparable à la convergence presque partout : aucune n'implique l'autre (3, 4). (iv) Multiplier deux distributions, ou diviser par une fonction qui s'annule (6, 9, 14). (v) Oublier la connexité dans le Théorème~6.4, ou les constantes perdues en dérivant (12, 13). (vi) Prendre une majoration de l'ordre pour une égalité (15).
\end{piege}

\end{document}
